\chapter{Digital Modulation and Optimal Detection}
\label{ch:modulation}

\begin{chapterintro}
Chapter~\ref{ch:baseband} sent bits as pulses on a wire. This chapter puts them on a carrier and
asks the two questions that define the whole discipline: \emph{which} set of waveforms should
represent the messages, and \emph{how} should a receiver decide which one was sent? The answers
come from one of the most beautiful ideas in engineering---that waveforms are vectors, noise is
a cloud around them, and optimal detection is a question of geometry. With that picture we
derive the maximum-likelihood receiver, compute the error rate of every classical constellation
(PAM, PSK, QAM, APSK, FSK), map the trade-off between power and bandwidth against Shannon's
limit, and meet the practical variants that real radios use: differential and noncoherent
detection, offset and $\pi/4$ schemes, MSK and GMSK, soft bit metrics for decoders, and the
measurements---EVM and MER---by which every transmitter on the market is judged. We finish at the
frontier: the constellations with deliberately unequal probabilities that now carry most of the
world's long-haul optical traffic.
\end{chapterintro}

\section{From bits to bandpass waveforms}
\label{sec:ch09:bandpass}

\subsection{What there is to modulate}

A sinusoid has three properties an engineer can vary: amplitude, phase and frequency.
Chapter~\ref{ch:analog} varied them continuously with a voice; a digital modulator varies them
in discrete steps, one step per \term{symbol}. Every $T$ seconds the modulator reads
$k=\log_2 M$ bits, chooses one of $M$ waveforms $s_1(t),\dots,s_M(t)$, and sends it. The
\term{symbol rate} is $R_s=1/T$ (in baud, after \'Emile Baudot) and the bit rate is
$R_b=kR_s$. The three classical families are named for what they vary:
\begin{itemize}
\item \term{amplitude-shift keying} (ASK) or, for a real signed alphabet, pulse-amplitude
modulation (PAM);
\item \term{phase-shift keying} (PSK), which varies only the phase;
\item \term{frequency-shift keying} (FSK), which varies the frequency;
\end{itemize}
and \term{quadrature amplitude modulation} (QAM) varies amplitude and phase together. In the
language of Chapter~\ref{ch:signals} every one of them is simply a rule for building the complex
envelope $\tilde s(t)$ in
\begin{equation}
s(t)=\re\{\tilde s(t)\,e^{j2\pi f_c t}\}=s_I(t)\cos 2\pi f_c t-s_Q(t)\sin 2\pi f_c t
\label{eq:ch09:bandpass}
\end{equation}
(equation~\eqref{eq:ch02:envelope}). The receiver's job is to recover $\tilde s(t)$ from a noisy,
distorted, frequency-shifted copy of $s(t)$, and then to decide which of the $M$ allowed
envelopes it most resembles.

\subsection{Linear modulation and the IQ modulator}

The workhorse of modern radio is \term{linear modulation}. Each symbol period the mapper turns
$k$ bits into one complex number $a_n=a_{I,n}+ja_{Q,n}$ drawn from a finite set, the
\term{constellation} or \emph{alphabet} $\mathcal A=\{\alpha_1,\dots,\alpha_M\}$, and the complex
envelope is a train of pulses weighted by those numbers:
\begin{equation}
\tilde s(t)=\sum_n a_n\,p(t-nT).
\label{eq:ch09:linmod}
\end{equation}
The pulse $p(t)$ is the root-raised-cosine or other Nyquist-compatible shape of
Chapter~\ref{ch:baseband}; the alphabet decides \emph{what} is sent, the pulse decides \emph{how
much spectrum} it takes. Because~\eqref{eq:ch09:linmod} is linear in the symbols, the spectrum of
$\tilde s(t)$ is just $|P(f)|^2$ scaled by the symbol variance (for zero-mean, uncorrelated
symbols), whatever the constellation. A 16-QAM and a QPSK signal with the same pulse and symbol
rate occupy exactly the same bandwidth; the 16-QAM signal simply carries twice as many bits in
it. That is the whole attraction of large constellations, and their price, as we shall see, is
power.

Figure~\ref{fig:ch09_iq} shows the hardware. The I and Q symbol streams are shaped by identical
filters, multiply a cosine and a negative sine from a single local oscillator, and are summed.
The receiver reverses the process: a quadrature downconverter produces $r_I(t)$ and $r_Q(t)$,
matched filters $p^*(-t)$ maximise the signal-to-noise ratio (Chapter~\ref{ch:baseband}), a
sampler takes one value per symbol at the right instant (Chapter~\ref{ch:sync}), and a
\term{detector} decides which constellation point each sample represents. In a software-defined
radio such as the B200 (Chapter~\ref{ch:sdr}) everything to the left of the DAC and right of the
ADC is arithmetic, but the structure is the same.

\begin{figure}[htbp]
\centering
\begin{tikzpicture}[font=\sffamily\scriptsize, >=Latex,
  blk/.style={draw=navy, fill=navy!6, thick, minimum height=7mm, minimum width=12mm, align=center},
  mix/.style={draw=navy, thick, circle, minimum size=5.5mm, inner sep=0pt},
  osc/.style={draw=accent, fill=accent!8, thick, minimum height=6mm, minimum width=12mm, align=center},
  lab/.style={color=navy!80},
  ttl/.style={font=\sffamily\scriptsize\bfseries, color=accent}]
% ---------------- transmitter
\node[ttl, anchor=west] at (-0.8,1.75) {TRANSMITTER};
\node[blk, minimum height=12mm] (map) at (0,0) {bits\\$\downarrow$\\mapper};
\node[blk] (pi) at (2.6,1.1) {$p(t)$};
\node[blk] (pq) at (2.6,-1.1) {$p(t)$};
\node[mix] (mi) at (4.9,1.1) {$\times$};
\node[mix] (mq) at (4.9,-1.1) {$\times$};
\node[osc] (lo) at (4.9,0) {LO $f_c$};
\node[mix] (sum) at (6.6,0) {$+$};
\node[blk] (pa) at (7.9,0) {PA};
\draw[->] (map.east) ++(0,0.3) -- ++(0.35,0) |- node[lab, above, pos=0.75] {$a_{I,n}$} (pi.west);
\draw[->] (map.east) ++(0,-0.3) -- ++(0.35,0) |- node[lab, below, pos=0.75] {$a_{Q,n}$} (pq.west);
\draw[->] (pi) -- node[lab, above] {$s_I(t)$} (mi);
\draw[->] (pq) -- node[lab, below] {$s_Q(t)$} (mq);
\draw[->] (lo) -- node[lab, right, pos=0.4] {$\cos 2\pi f_ct$} (mi);
\draw[->] (lo) -- node[lab, right, pos=0.4] {$-\sin 2\pi f_ct$} (mq);
\draw[->] (mi) -| (sum);
\draw[->] (mq) -| (sum);
\draw[->] (sum) -- (pa);
\draw[->] (pa.east) -- ++(0.5,0) node[right, lab] {$s(t)$};
% ---------------- receiver
\begin{scope}[yshift=-4.3cm]
\node[ttl, anchor=west] at (-0.8,1.75) {RECEIVER};
\node[blk] (lna) at (0,0) {LNA};
\node[mix] (ri) at (1.9,1.1) {$\times$};
\node[mix] (rq) at (1.9,-1.1) {$\times$};
\node[osc] (lo2) at (1.9,0) {LO $f_c$};
\node[blk] (mfi) at (4.4,1.1) {$p^*(-t)$};
\node[blk] (mfq) at (4.4,-1.1) {$p^*(-t)$};
\node[blk, minimum height=30mm, minimum width=17mm] (det) at (7.0,0) {sample\\at $t=nT$\\[4pt]nearest-\\point\\detector};
\draw[<-] (lna.west) -- ++(-0.5,0) node[left, lab] {$r(t)$};
\draw[->] (lna.east) -- ++(0.35,0) |- (ri.west);
\draw[->] (lna.east) -- ++(0.35,0) |- (rq.west);
\draw[->] (lo2) -- node[lab, right, pos=0.4] {$\cos$} (ri);
\draw[->] (lo2) -- node[lab, right, pos=0.4] {$-\sin$} (rq);
\draw[->] (ri) -- node[lab, above] {$r_I(t)$} (mfi);
\draw[->] (rq) -- node[lab, below] {$r_Q(t)$} (mfq);
\draw[->] (mfi.east) -- (det.west |- mfi.east);
\draw[->] (mfq.east) -- (det.west |- mfq.east);
\draw[->] (det.east) -- ++(0.5,0) node[right, lab] {$\hat a_n\to$ bits};
\end{scope}
\end{tikzpicture}
\caption{The IQ modulator and demodulator of linear digital modulation. Two real pulse trains
ride on quadrature carriers from one oscillator; the receiver separates them, matched-filters,
samples once per symbol and decides. Synchronization (Chapter~\ref{ch:sync}) supplies the
carrier phase and sampling instant this diagram takes for granted.}
\label{fig:ch09_iq}
\end{figure}

\subsection{Nonlinear modulation}

Not every useful scheme fits~\eqref{eq:ch09:linmod}. In FSK and its refined descendants---the
continuous-phase modulations MSK and GMSK of Section~\ref{sec:ch09:constenv}---the bits steer the
\emph{frequency}, and the envelope is $e^{j\phi(t)}$ with a phase that remembers every past
bit. These signals have constant amplitude, which a power amplifier loves, but they do not have
the tidy spectral properties of linear modulation, and their optimal receivers need memory. We
develop the geometric theory first for arbitrary signal sets, so that it covers linear and
nonlinear schemes alike, and return to the nonlinear family once the tools are in hand.

\section{Signal space}
\label{sec:ch09:signalspace}

\subsection{Waveforms are vectors}

The central trick of detection theory is to stop thinking about waveforms. Signals of finite
energy on an interval $[0,T)$ form a vector space with inner product and norm
\begin{equation}
\langle x,y\rangle=\int_0^T x(t)\,y^*(t)\,dt,\qquad \|x\|^2=\langle x,x\rangle=E_x .
\label{eq:ch09:inner}
\end{equation}
Suppose the $M$ signals of a modulation all lie in an $N$-dimensional subspace ($N\le M$)
spanned by \term{orthonormal basis functions} $\phi_1(t),\dots,\phi_N(t)$, meaning
$\langle\phi_m,\phi_n\rangle=\delta_{mn}$. Then each signal is a weighted sum of basis
functions,
\begin{equation}
s_i(t)=\sum_{n=1}^{N}s_{in}\,\phi_n(t),\qquad s_{in}=\langle s_i,\phi_n\rangle,
\label{eq:ch09:expansion}
\end{equation}
and is completely described by its coefficient vector $\vect s_i=(s_{i1},\dots,s_{iN})$. The
map from waveforms to vectors is an \emph{isometry}: energies and distances are preserved,
\begin{equation}
E_i=\int|s_i(t)|^2dt=\|\vect s_i\|^2,\qquad
\int|s_i(t)-s_j(t)|^2dt=\|\vect s_i-\vect s_j\|^2 .
\label{eq:ch09:isometry}
\end{equation}
This is Parseval's theorem in disguise, and it means that every question about the energy or
distinguishability of a set of waveforms becomes a question of plane (or $N$-dimensional)
geometry.

Figure~\ref{fig:ch09_signal_space} makes the idea concrete for QPSK. The four waveforms are
cosine bursts with phases $45^\circ$, $135^\circ$, $225^\circ$ and $315^\circ$; all four are
combinations of just two basis functions, a cosine and a negative sine of energy one. In signal
space the modulation is four points on a circle of radius $\sqrt{E_s}$. The carrier frequency,
the waveform shapes and the duration have vanished; what remains is the only thing that matters
to a detector.

\mdcfig{ch09_signal_space}{From waveforms to points. Left: the four QPSK waveforms over one
symbol. Centre: the two orthonormal basis functions that span them. Right: the signal-space
representation, four points on a circle of radius $\sqrt{E_s}$. Distances between points equal
the root energy of the difference between waveforms.}

For a linear modulation with a pulse $p(t)$ of unit energy, the basis functions are
\begin{equation}
\phi_1(t)=\sqrt2\,p(t)\cos 2\pi f_ct,\qquad \phi_2(t)=-\sqrt2\,p(t)\sin 2\pi f_ct
\end{equation}
(orthogonal when $f_c$ is much
larger than the pulse bandwidth), and the signal-space point of a symbol is simply the complex
number $a_n$ read as a two-vector. This is why engineers draw constellations in the complex
plane and why this book uses the words ``constellation point'', ``symbol'' and ``signal vector''
interchangeably. FSK is different: $M$ tones with suitable spacing are mutually orthogonal, so
$M$-FSK needs $N=M$ dimensions, a fact whose consequences fill Section~\ref{sec:ch09:orthogonal}.

\subsection{Finding a basis: Gram--Schmidt}

Given any finite set of signals, the \term{Gram--Schmidt procedure} builds an orthonormal basis
for them. Take the first signal and normalise it: $\phi_1=s_1/\|s_1\|$. For each subsequent
signal, subtract its projections on the basis functions found so far and normalise what is
left, if anything:
\begin{equation}
g_i(t)=s_i(t)-\sum_{n<i}\langle s_i,\phi_n\rangle\phi_n(t),\qquad
\phi_i(t)=\frac{g_i(t)}{\|g_i\|}\quad\text{if }\|g_i\|>0 .
\label{eq:ch09:gs}
\end{equation}
If the residual $g_i$ is zero, $s_i$ already lies in the span and adds no dimension. The
procedure always terminates with $N\le M$ basis functions. The basis is not unique---a different
ordering or any rotation gives another---but the \emph{geometry} (all lengths and angles) is the
same in every basis, and so is every error probability.

\begin{worked}[Gram--Schmidt on three rectangular signals]
Let $T=2$ and define $s_1(t)=1$ on $[0,2)$; $s_2(t)=2$ on $[0,1)$ and $0$ on $[1,2)$;
$s_3(t)=+1$ on $[0,1)$ and $-1$ on $[1,2)$.

\emph{Step 1.} $\|s_1\|^2=2$, so $\phi_1=s_1/\sqrt2$ and $\vect s_1=(\sqrt2,0)$.

\emph{Step 2.} $\langle s_2,\phi_1\rangle=\int_0^1 2\cdot\tfrac1{\sqrt2}\,dt=\sqrt2$. The
residual $g_2=s_2-\sqrt2\,\phi_1=s_2-s_1$ equals $+1$ on $[0,1)$ and $-1$ on $[1,2)$, with
energy 2, so $\phi_2=g_2/\sqrt2$ and $\vect s_2=(\sqrt2,\sqrt2)$.

\emph{Step 3.} $\langle s_3,\phi_1\rangle=0$ and $\langle s_3,\phi_2\rangle=\sqrt2$. The residual
is zero: $s_3=s_2-s_1$ adds no dimension, and $\vect s_3=(0,\sqrt2)$.

Three signals live in a plane. The distances from the vectors are
$\|\vect s_1-\vect s_2\|=\sqrt2$, $\|\vect s_1-\vect s_3\|=2$ and
$\|\vect s_2-\vect s_3\|=\sqrt2$. Check one directly: $s_1-s_3$ is $0$ on $[0,1)$ and $2$ on
$[1,2)$, energy 4, distance 2. The receiver for this set needs two correlators, not three.
\end{worked}

\subsection{The vector channel}

Now send $s_i(t)$ through the additive white Gaussian noise channel of Chapter~\ref{ch:noise},
\begin{equation}
r(t)=s_i(t)+w(t),
\end{equation}
where $w(t)$ is white with two-sided power spectral density $N_0/2$ (for real passband signals;
for complex envelopes the noise is circularly symmetric with variance $N_0$ per complex sample
at the output of a unit-energy matched filter). Project the received waveform onto each basis
function:
\begin{equation}
r_n=\langle r,\phi_n\rangle=s_{in}+w_n,\qquad n=1,\dots,N .
\label{eq:ch09:project}
\end{equation}
Because the $\phi_n$ are orthonormal and the noise is white, the noise coordinates $w_n$ are
independent Gaussian random variables, each of variance $N_0/2$:
\begin{equation}
\E[w_mw_n]=\int\!\!\int\phi_m(t)\,\frac{N_0}{2}\delta(t-u)\,\phi_n(u)\,dt\,du=\frac{N_0}{2}\delta_{mn}.
\end{equation}
The whole waveform channel has become a \term{vector channel},
\begin{equation}
\vect r=\vect s_i+\vect w,\qquad \vect w\sim\mathcal N\!\left(\vect 0,\tfrac{N_0}{2}\vect I_N\right),
\label{eq:ch09:vecchan}
\end{equation}
and the received vector is a sample from a spherical Gaussian cloud centred on the transmitted
point. The spherical shape is the key: white noise has no preferred direction, so it disturbs
every coordinate equally, and a rotation of the constellation changes nothing.

\subsection{Sufficient statistics and the theorem of irrelevance}

Something seems to have been thrown away. The received waveform $r(t)$ is infinite-dimensional;
its projection $\vect r$ has only $N$ components. The part left over,
\begin{equation}
r'(t)=r(t)-\sum_{n=1}^N r_n\phi_n(t)=w(t)-\sum_n w_n\phi_n(t),
\end{equation}
contains no signal at all---only the component of the noise orthogonal to the signal space. And
because Gaussian random variables that are uncorrelated are independent, $r'(t)$ is independent
of the $w_n$ and hence of everything the receiver cares about. It carries no information about
which signal was sent, and no information about the noise in the signal space either. This is
the \term{theorem of irrelevance} (Wozencraft and Jacobs, 1965): for signals in white Gaussian
noise, the projections $\vect r$ are a \term{sufficient statistic}; the optimal receiver may
discard everything else without loss.

\begin{keyidea}[The dimensionality of detection]
A receiver for $M$ signals in white Gaussian noise never needs to process more than $N\le M$
numbers per symbol: the correlations of the received waveform with an orthonormal basis for the
signals. For every linear modulation, however large the constellation, $N=2$---one I sample and
one Q sample from a matched filter. Everything else in $r(t)$ is irrelevant noise.
\end{keyidea}

The theorem has conditions, and practising engineers must know where they break. The noise must
be white (or whitened first, as the equalisers of Chapter~\ref{ch:equalization} do) and
Gaussian; interference from other users is neither, and a receiver that exploits its structure
(Chapter~\ref{ch:mimo}) can beat one that treats it as noise. The signals must be known exactly;
an unknown carrier phase, a timing error or a frequency offset makes the true signals differ from
the ones the correlators expect, and the out-of-subspace residual is then no longer irrelevant.
And the argument is per symbol; when pulses overlap (intersymbol interference), the sufficient
statistic is the whole sequence of matched-filter outputs, which is where the sequence detectors
of Chapter~\ref{ch:equalization} begin.

\section{Optimal detection}
\label{sec:ch09:detection}

\subsection{MAP and ML decision rules}

The receiver observes $\vect r$ and must choose $\hat\imath$. The rule that minimises the
probability of error is the one that picks the most probable signal given the observation, the
\term{maximum a posteriori} (MAP) rule:
\begin{equation}
\hat\imath=\arg\max_i\,P(\vect s_i\mid\vect r)=\arg\max_i\,p_i\,f(\vect r\mid\vect s_i),
\label{eq:ch09:map}
\end{equation}
where $p_i$ is the prior probability of symbol $i$ and the second form follows from Bayes' rule
(the denominator $f(\vect r)$ is common to all $i$). The proof is one line: the probability of a
correct decision is $\int P(\vect s_{\hat\imath(\vect r)}\mid\vect r)f(\vect r)\,d\vect r$, which is
maximised pointwise by choosing, for each $\vect r$, the $i$ with the largest posterior. When the
priors are equal---as they are for scrambled, compressed or encrypted data---the MAP rule reduces
to the \term{maximum-likelihood} (ML) rule, $\hat\imath=\arg\max_i f(\vect r\mid\vect s_i)$.

For the Gaussian vector channel~\eqref{eq:ch09:vecchan} the likelihood is
\begin{equation}
f(\vect r\mid\vect s_i)=(\pi N_0)^{-N/2}\exp\!\left(-\frac{\|\vect r-\vect s_i\|^2}{N_0}\right),
\end{equation}
so taking logarithms,
\begin{equation}
\hat\imath_{\text{MAP}}=\arg\min_i\Big[\|\vect r-\vect s_i\|^2-N_0\ln p_i\Big],
\qquad
\hat\imath_{\text{ML}}=\arg\min_i\|\vect r-\vect s_i\|^2 .
\label{eq:ch09:mindist}
\end{equation}
The ML detector in white Gaussian noise is a \term{minimum-distance detector}: choose the
constellation point nearest the received vector. The MAP detector adds a bias that favours
likely symbols, weighted by $N_0$---negligible at high SNR, important at low SNR and whenever
priors are strongly unequal, as they are in the shaped constellations of
Section~\ref{sec:ch09:shaping} and in iterative decoders that feed prior information back to the
demapper.

\subsection{Decision regions}

The minimum-distance rule partitions signal space into \term{decision regions}: the region
$\mathcal D_i$ is the set of points closer to $\vect s_i$ than to any other signal point. The
boundary between two neighbours is the perpendicular bisector of the segment joining them, and
the regions are the \term{Voronoi cells} of the constellation. For square QAM the cells are
squares (and half-open strips at the edges), so the two-dimensional decision separates into two
independent one-dimensional slicers, one for I and one for Q---the reason square QAM is so cheap
to demodulate. For PSK the cells are pie wedges, and the detector need only measure the phase.
For APSK, with its concentric rings, the cells mix wedges and annuli
(Figure~\ref{fig:ch09_decision_regions}).

\mdcfig{ch09_decision_regions}{Decision regions of the minimum-distance detector. Left: square
16-QAM, whose Voronoi cells are rectangles, with received samples at $E_s/N_0=16$\,dB. Centre:
16-APSK (4+12), where the inner ring's cells are bounded by a polygon and the outer cells are
wedges. Right: with unequal priors the MAP boundary between two antipodal points moves towards
the less likely one; here $p_0=0.8$ shifts it from $0$ to $+0.44$ at $\sigma=0.8$.}

With unequal priors the boundaries shift. For two antipodal points $\pm A$ in Gaussian noise of
variance $\sigma^2$, equating the two MAP metrics gives the threshold
\begin{equation}
\tau=\frac{\sigma^2}{2A}\ln\frac{p_0}{p_1},
\label{eq:ch09:mapthresh}
\end{equation}
where $p_0$ is the prior of $-A$. The boundary slides away from the more probable symbol, giving
it a larger region. In the example of the figure ($p_0=0.8$, $A=1$, $\sigma=0.8$) the threshold
moves to $0.44$, and the error probability drops from $\Q(1.25)=0.106$ with the ML threshold to
$0.8\,\Q(1.80)+0.2\,\Q(0.70)=0.077$---a real gain, bought purely by using what the receiver
knew about the source.

\subsection{The correlation receiver}

Expanding the squared distance in~\eqref{eq:ch09:mindist},
\begin{equation}
\|\vect r-\vect s_i\|^2=\|\vect r\|^2-2\re\langle\vect r,\vect s_i\rangle+E_i ,
\end{equation}
and dropping the term common to all hypotheses gives an equivalent rule,
\begin{equation}
\hat\imath=\arg\max_i\left[\re\langle\vect r,\vect s_i\rangle-\frac{E_i}{2}+\frac{N_0}{2}\ln p_i\right].
\label{eq:ch09:corr}
\end{equation}
The receiver correlates the input against each candidate, subtracts half its energy (so that a
high-energy signal does not win merely by being large) and, for MAP, adds a prior bias. For
equal-energy sets such as PSK and FSK with equal priors, the rule is simply \emph{pick the largest
correlation}. The correlations can be computed in signal space ($N$ correlators followed by a
small matrix multiply) or directly against the waveforms ($M$ correlators); either way, each
correlator can be replaced by a filter matched to the basis function and sampled at $t=T$, the
matched filter of Chapter~\ref{ch:baseband}. Figure~\ref{fig:ch09_correlator} shows the
structure.

\begin{figure}[htbp]
\centering
\begin{tikzpicture}[font=\sffamily\scriptsize, >=Latex,
  blk/.style={draw=navy, fill=navy!6, thick, minimum height=6.5mm, minimum width=11mm, align=center},
  mix/.style={draw=navy, thick, circle, minimum size=5mm, inner sep=0pt},
  lab/.style={color=navy!80}]
\node (in) at (0,0) {$r(t)$};
\foreach \y/\n/\lbl in {1.1/1/{\phi_1(t)}, 0/2/{\phi_2(t)}, -1.3/N/{\phi_N(t)}}{
  \node[mix] (m\n) at (1.6,\y) {$\times$};
  \node[blk, right=5mm of m\n] (int\n) {$\int_0^T\!(\cdot)\,dt$};
  \node[above=0.5mm of m\n, lab] {$\lbl$};
  \draw[->] (in.east) -- ++(0.3,0) |- (m\n.west);
  \draw[->] (m\n) -- (int\n);
}
\node at (1.6,-0.6) {$\vdots$};
\node[blk, minimum height=30mm, minimum width=24mm, right=9mm of int2] (metric)
  {compute for\\each $i=1..M$:\\[3pt]$\re\langle\vect r,\vect s_i\rangle$\\$-\;E_i/2$\\$+\;\tfrac{N_0}{2}\ln p_i$};
\foreach \n in {1,2,N}{\draw[->] (int\n.east) -- node[lab, above] {$r_{\n}$} (metric.west |- int\n.east);}
\node[blk, right=6mm of metric, fill=accent!8, draw=accent] (am) {choose\\largest};
\draw[->] (metric) -- (am);
\draw[->] (am.east) -- ++(5mm,0) node[right, lab] {$\hat\imath$};
\draw[decorate, decoration={brace, amplitude=4pt, mirror}, navy!60] (1.2,-1.8) -- (3.9,-1.8)
  node[midway, below=4pt, lab] {sufficient statistic $\vect r$ ($N$ numbers)};
\end{tikzpicture}
\caption{The optimal (MAP) receiver for $M$ signals in white Gaussian noise. $N$ correlators (or
matched filters sampled at $T$) project the input onto the signal space; the metric block
evaluates~\eqref{eq:ch09:corr}. For linear modulation $N=2$ and the metric block reduces to a
nearest-point slicer.}
\label{fig:ch09_correlator}
\end{figure}

\begin{history}[Kotelnikov, Shannon and the geometry of signals]
The geometric view of signals has two independent origins. In his 1947 doctoral thesis, published
in English in 1959 as \emph{The Theory of Optimum Noise Immunity}, Vladimir Kotelnikov represented
signals as vectors, derived the optimal receiver for signals in Gaussian noise, and computed the
error probabilities of amplitude, frequency and phase keying---work essentially unknown in the
West for a decade. Claude Shannon's 1949 paper ``Communication in the presence of noise'' used
the same geometry to prove the capacity formula, treating a band-limited signal of duration $T$
as a point in a space of $2WT$ dimensions and noise as a small sphere around it. It was John
Wozencraft and Irwin Jacobs's 1965 textbook \emph{Principles of Communication Engineering} that
made signal space the standard language of the field; its theorem of irrelevance is still quoted
in their words. Jacobs went on to co-found Linkabit and then Qualcomm, where the geometry became
CDMA handsets by the hundred million.
\end{history}

\section{Error probability}
\label{sec:ch09:errorprob}

\subsection{Two signals: the pairwise error probability}

Everything starts with two points. If $\vect s_i$ is sent and only $\vect s_j$ competes, the ML
detector errs when the noise carries $\vect r$ across the bisector between them. Only the
component of the noise along the line joining the points matters---the perpendicular components
move $\vect r$ parallel to the boundary---and that component is a one-dimensional Gaussian of
variance $N_0/2$. The boundary is $d_{ij}/2$ away, where $d_{ij}=\|\vect s_i-\vect s_j\|$, so
\begin{equation}
P(\vect s_i\to\vect s_j)=\Q\!\left(\frac{d_{ij}/2}{\sqrt{N_0/2}}\right)=\Q\!\left(\sqrt{\frac{d_{ij}^2}{2N_0}}\right),
\label{eq:ch09:pairwise}
\end{equation}
with $\Q(\cdot)$ the Gaussian tail function of equation~\eqref{eq:ch03:q}. This
\term{pairwise error probability} depends on nothing but the distance between the two signals
and the noise density: not their shapes, not their individual energies, not the carrier. For
antipodal signals $\pm\sqrt{E_b}$, $d^2=4E_b$ and $P_b=\Q(\sqrt{2E_b/N_0})$; for orthogonal
signals of energy $E_b$, $d^2=2E_b$ and $P_b=\Q(\sqrt{E_b/N_0})$, the 3\,dB penalty found in
Chapter~\ref{ch:baseband}. Distance is the currency of detection.

\subsection{Many signals: the union bound}

With $M$ signals the exact error probability is an integral of the Gaussian cloud over the
outside of a polygonal (or polyhedral) region, rarely available in closed form. But an error
means the noise has carried $\vect r$ closer to \emph{some} other point, and the probability of a
union of events never exceeds the sum of their probabilities. Hence the \term{union bound}:
\begin{equation}
P_s\le\frac1M\sum_{i=1}^M\sum_{j\ne i}\Q\!\left(\sqrt{\frac{d_{ij}^2}{2N_0}}\right)
\label{eq:ch09:union}
\end{equation}
for equiprobable signals. It overcounts regions where several pairwise error events overlap,
which matters only when errors are common. At the error rates engineers care about, the bound is
dominated by the smallest distances, because $\Q$ falls so steeply: a competitor at
$\sqrt2\,d_{\min}$ contributes roughly a factor $e^{-d_{\min}^2/(4N_0)}$ less than a nearest
neighbour, which at $10^{-5}$ error rates is a factor of several thousand. Keeping only the nearest
neighbours gives the \term{nearest-neighbour approximation}
\begin{equation}
P_s\approx\bar N_{\min}\,\Q\!\left(\sqrt{\frac{d_{\min}^2}{2N_0}}\right),
\label{eq:ch09:nn}
\end{equation}
where $d_{\min}$ is the \term{minimum distance} of the constellation and $\bar N_{\min}$ the
average number of neighbours at that distance. Figure~\ref{fig:ch09_union_bound} compares the
exact result, the full bound and the approximation for 8-PSK and 16-QAM: below a symbol error
rate of about $10^{-2}$ they are indistinguishable.

\mdcfig{ch09_union_bound}{Exact symbol error probability, the union bound~\eqref{eq:ch09:union}
and the nearest-neighbour approximation~\eqref{eq:ch09:nn}, with Monte Carlo points (400\,000
symbols each), for 8-PSK and 16-QAM. The bounds are loose only where errors are too frequent to be
interesting.}

\begin{keyidea}[Two numbers characterise a constellation]
At high SNR the error probability of any constellation in white Gaussian noise is set by two
numbers: the normalised minimum distance $d_{\min}^2/E_s$, which sets the position of the
waterfall, and the neighbour count $\bar N_{\min}$, which shifts it slightly. When constellations
of different sizes are compared at equal energy per \emph{bit}, $d_{\min}^2/E_b$ is the figure of
merit. Every decibel lost in $d_{\min}^2/E_b$ is a decibel of transmit power.
\end{keyidea}

\subsection{Bit errors and symbol errors}

The union bound speaks of symbol errors, but users and decoders count bits. The relationship
depends on the \term{labelling}---which $k$-bit pattern is assigned to which point. Almost every
symbol error at useful SNR is to a nearest neighbour, so if neighbours differ in only one bit,
a symbol error costs one bit error out of $k$:
\begin{equation}
P_b\approx\frac{P_s}{k}\qquad\text{(Gray labelling)} .
\label{eq:ch09:graybits}
\end{equation}
A labelling in which every pair of nearest neighbours differs in exactly one bit is a
\term{Gray labelling}, after Frank Gray of Bell Labs, whose patent on the reflected binary
code (filed 1947, granted 1953) was motivated by an electron-beam tube for PCM encoding.
For PSK, labelling the points around the circle with
the reflected binary sequence 000, 001, 011, 010, 110, 111, 101, 100 achieves it; for square QAM, Gray
labelling each axis independently does. Cross constellations (32 and 128 points) cannot be
perfectly Gray-labelled, and the best labellings leave a few neighbour pairs differing in two
bits, a penalty of a small fraction of a decibel. Figure~\ref{fig:ch09_gray_mapping} compares
Gray and natural binary labels on 16-QAM. The difference is smaller than folklore suggests---about
a third more bit errors, a couple of tenths of a decibel---but it is free to avoid, and it
matters more in coded systems, where the labelling also sets the reliability of the individual
bits the decoder sees (Section~\ref{sec:ch09:soft}).

\mdcfig{ch09_gray_mapping}{Gray versus natural binary labels on 16-QAM. With Gray labels,
adjacent points differ in one bit; with natural labels some neighbours differ in two (arrows).
The bit error rate with Gray labels follows $P_s/k$ closely; natural labels cost about a third
more bit errors.}

\section{The classical constellations}
\label{sec:ch09:constellations}

Figure~\ref{fig:ch09_constellations} is a gallery of the constellations that carry most of the
world's digital radio traffic, each normalised to unit average energy, with its normalised
minimum distance and its peak-to-average power ratio. We derive their error rates in turn.

\mdcfig{ch09_constellations}{A constellation gallery at unit average energy. Labels on 8-PSK and
16-QAM are Gray. The figures beneath each panel give the squared minimum distance normalised to
the average symbol energy and the ratio of peak to average constellation power; APSK buys a
lower peak-to-average ratio than QAM at some cost in distance.}

\subsection{\texorpdfstring{$M$}{M}-ary PAM}

The one-dimensional constellation of Chapter~\ref{ch:baseband} places $M$ levels at
$\pm A,\pm3A,\dots,\pm(M-1)A$, with $d_{\min}=2A$ and average energy
$E_s=A^2(M^2-1)/3$. The two outer points have one neighbour each and the $M-2$ inner points two,
so $\bar N_{\min}=2(M-1)/M$, and here the nearest-neighbour expression is exact:
\begin{equation}
P_s^{\text{PAM}}=2\left(1-\frac1M\right)\Q\!\left(\sqrt{\frac{6}{M^2-1}\,\frac{E_s}{N_0}}\right).
\label{eq:ch09:pam}
\end{equation}
Each doubling of $M$ adds one bit per symbol and multiplies $M^2-1$ by roughly four: 6\,dB more
energy per symbol for the same error rate. That exchange rate---6\,dB per bit per dimension---is
the steep price of amplitude levels, and it recurs in QAM, in the quantization rule of
Chapter~\ref{ch:pcm} and in Shannon's formula at high SNR.

\subsection{\texorpdfstring{$M$}{M}-ary PSK}

Phase-shift keying places $M$ points on a circle: $a=\sqrt{E_s}\,e^{j(2\pi m/M+\theta_0)}$. All
points have the same energy---a constant envelope at the symbol instants, though not between them
once the pulses are shaped (Section~\ref{sec:ch09:constenv}). Neighbours are separated by the
chord $d_{\min}=2\sqrt{E_s}\sin(\pi/M)$, each point has two, and~\eqref{eq:ch09:nn} gives
\begin{equation}
P_s^{\text{PSK}}\approx2\,\Q\!\left(\sqrt{\frac{2E_s}{N_0}}\sin\frac\pi M\right) .
\label{eq:ch09:psk}
\end{equation}
An exact result, easy to integrate numerically, is Craig's formula (1991):
\begin{equation}
P_s^{\text{PSK}}=\frac1\pi\int_0^{(M-1)\pi/M}\exp\!\left(-\frac{E_s}{N_0}\,\frac{\sin^2(\pi/M)}{\sin^2\theta}\right)d\theta .
\label{eq:ch09:craig}
\end{equation}
Two special cases deserve their own names. \term{BPSK} ($M=2$) is antipodal signalling,
$P_b=\Q(\sqrt{2E_b/N_0})$. \term{QPSK} ($M=4$) is two BPSK signals in quadrature, one on I and one
on Q; with Gray labels each bit sees an independent BPSK channel with the same $E_b$, so
\begin{equation}
P_b^{\text{QPSK}}=\Q\!\left(\sqrt{2E_b/N_0}\right),
\end{equation}
identical to BPSK. QPSK doubles the bit rate in the same bandwidth for no loss in power
efficiency---the nearest thing to a free lunch in this subject, and the reason QPSK is the default
robust mode of nearly every modern system. Beyond four points, PSK becomes expensive. For large
$M$, $\sin(\pi/M)\approx\pi/M$ and each doubling of $M$ costs 6\,dB of $E_s/N_0$, the same as
PAM, because the points are strung along a one-dimensional circle. 8-PSK needs
$E_b/N_0\approx13.0$\,dB for $P_b=10^{-5}$, 3.4\,dB more than QPSK; 16-PSK needs 17.4\,dB.

\subsection{Square QAM}

Quadrature amplitude modulation uses both dimensions for amplitude. Square $M$-QAM with $M=L^2$
is two independent $L$-PAM signals, one on I and one on Q, each carrying half the energy. From the
PAM result, the probability that one axis errs is
\begin{equation}
P_L=2\left(1-\frac1{\sqrt M}\right)\Q\!\left(\sqrt{\frac{3}{M-1}\,\frac{E_s}{N_0}}\right),
\end{equation}
and a symbol is correct only if both axes are, so exactly
\begin{equation}
P_s^{\text{QAM}}=1-(1-P_L)^2=2P_L-P_L^2 .
\label{eq:ch09:qam}
\end{equation}
The normalised minimum distance is $d_{\min}^2=6E_s/(M-1)$. With Gray labels on each axis,
\begin{equation}
P_b^{\text{QAM}}\approx\frac4{\log_2M}\left(1-\frac1{\sqrt M}\right)
\Q\!\left(\sqrt{\frac{3\log_2M}{M-1}\,\frac{E_b}{N_0}}\right).
\label{eq:ch09:qamber}
\end{equation}
Compare PSK and QAM at 16 points: $d_{\min}^2/E_s$ is $4\sin^2(\pi/16)=0.152$ for 16-PSK and
$6/15=0.400$ for 16-QAM, a 4.2\,dB advantage for QAM, because QAM spreads its points over the
disc rather than confining them to its rim. At every size above eight points, QAM beats PSK in
white Gaussian noise; that is why PSK stops at 8 (in EDGE and DVB-S2) and QAM goes to 4096.

Figure~\ref{fig:ch09_ber_families} collects the bit error rates. Each factor of four in QAM size
(two more bits per symbol) costs about 6\,dB of $E_s/N_0$ and about 4.4 to 5.2\,dB of $E_b/N_0$ at a
given error rate. Small constellations are power-efficient; large ones are
bandwidth-efficient. Section~\ref{sec:ch09:tradeoff} makes that trade precise.

\begin{worked}[How much $E_b/N_0$ does uncoded 64-QAM need?]
For $M=64$, $k=6$, equation~\eqref{eq:ch09:qamber} reads
\[
P_b\approx\tfrac46\left(1-\tfrac18\right)\Q\!\left(\sqrt{\tfrac{18}{63}\,E_b/N_0}\right)
=0.583\,\Q\!\left(\sqrt{0.2857\,E_b/N_0}\right).
\]
For $P_b=10^{-5}$ we need $\Q(x)=1.71\times10^{-5}$, so $x\approx4.14$ (from a table, or
\texttt{np.sqrt(2)*erfcinv(2*1.71e-5)} in Python). Then $E_b/N_0=4.14^2/0.2857=60.0$, which is
\textbf{17.8\,dB}, and $E_s/N_0=17.8+10\log_{10}6=$ \textbf{25.6\,dB}.

For comparison, QPSK needs 9.6\,dB. Tripling the spectral efficiency from 2 to 6 bits per symbol
costs 8.2\,dB more energy per bit; at a fixed \emph{symbol} rate, where the bit rate triples too,
the received power must rise by $8.2+4.8=13$\,dB. Forward error correction moves these numbers a
long way: with a rate-5/6 LDPC code, 64-QAM delivers 5 information bits per symbol at an
$E_s/N_0$ of roughly 17 to 18\,dB, within about 3\,dB of the Shannon limit of 14.9\,dB for that
rate (Chapter~\ref{ch:moderncodes}).
\end{worked}

\mdcfig{ch09_ber_families}{Bit error probability of Gray-labelled $M$-PSK (left) and square
$M$-QAM (right) against $E_b/N_0$. Lines are exact (PSK, via~\eqref{eq:ch09:craig} divided by
$k$) or the nearest-neighbour approximation~\eqref{eq:ch09:qamber}; circles are simulations with
\texttt{commlib}. The dotted line marks $10^{-5}$.}


\subsection{Cross and rectangular QAM}

When $k$ is odd, $M=2^k$ is not a perfect square. A rectangular $4\times8$ grid works but wastes
energy on its long axis. The \term{cross constellation} does better: take a $6\times6$ grid and
remove the four corners (32-QAM), or a $12\times12$ grid and remove $2\times2$ blocks from each
corner (128-QAM). The corners are the most expensive points in energy, so removing them lowers
the average energy for the same minimum distance. For 32 points, $d_{\min}^2/E_s$ is $0.200$ for
the cross against $0.154$ for the $4\times8$ rectangle, a 1.1\,dB advantage. Cross constellations
were workhorses of the voiceband-modem era (the history box below) and 32- and 128-QAM remain
options in DVB-C cable television and in microwave backhaul. Their Gray labelling is imperfect, so
they have largely given way to square QAM combined with the fine rate granularity of modern codes.

\subsection{APSK}

Satellites use \term{amplitude and phase-shift keying} (APSK): points on a few concentric rings,
such as the 4+12 arrangement of 16-APSK and the 4+12+16 arrangement of 32-APSK in DVB-S2 (ETSI
EN~302~307-1). The motive is the amplifier. A satellite travelling-wave-tube amplifier is most
efficient near saturation, where it compresses amplitude and rotates phase by an amount that
depends on the input amplitude (AM/AM and AM/PM conversion, Chapter~\ref{ch:sdr}). Such
distortion is hard to correct across many amplitude levels (16-QAM has three distinct amplitudes,
64-QAM nine) but easy across two or three rings: the ground station simply predistorts each ring's
radius and phase. APSK also has a lower peak-to-average ratio than QAM of the same size
(Figure~\ref{fig:ch09_constellations}), so the amplifier can run closer to saturation. The cost is
distance: 16-APSK with ring ratio $\gamma=R_2/R_1=2.85$ has $d_{\min}^2/E_s=0.315$ against
16-QAM's 0.400, about 1\,dB worse on a perfectly linear channel. On a real satellite the
nonlinearity erases that deficit and more. DVB-S2 specifies the ring ratio as a function of code
rate---from 3.15 at rate 2/3 down to about 2.6 at rate 9/10---and DVB-S2X (EN~302~307-2) extends
the family to 64-, 128- and 256-APSK, used for professional contribution links and high-throughput
backhaul (Chapter~\ref{ch:satellite}).

\begin{inpractice}[High-order QAM in the wild]
Constellation sizes have climbed steadily as radios have become cleaner and SNRs higher:
\begin{itemize}
\item \textbf{Wi-Fi}: 802.11a/g stopped at 64-QAM, 802.11ac (Wi-Fi~5) added 256-QAM, 802.11ax
(Wi-Fi~6) 1024-QAM and 802.11be (Wi-Fi~7) 4096-QAM. The larger constellations are usable only
close to the access point, with a strong signal and excellent transmitter and receiver EVM
(Table~\ref{tab:ch09:wifimcs}).
\item \textbf{Cellular}: LTE began with 64-QAM and added 256-QAM in the downlink (Release 12);
5G NR supports up to 256-QAM in both directions and 1024-QAM in the downlink from Release 17
(3GPP TS 38.211, TS 38.214).
\item \textbf{Cable}: DOCSIS 3.0 used 64- and 256-QAM downstream channels 6 or 8\,MHz wide.
DOCSIS 3.1 moved to OFDM with subcarrier constellations up to 4096-QAM downstream (8192 and 16384
are defined as options) and up to 1024-QAM upstream---practical because a well-maintained coaxial
plant offers SNRs around 40\,dB.
\item \textbf{Microwave backhaul}: point-to-point radios with fixed high-gain antennas and adaptive
modulation run from QPSK in heavy rain to 4096-QAM in clear sky, switching without dropping
traffic.
\item \textbf{Optical}: coherent transceivers use QPSK to 64-QAM on each polarisation, increasingly
with probabilistic shaping (Section~\ref{sec:ch09:shaping}).
\end{itemize}
The pattern is universal: the constellation is chosen per link, often per frame, from the measured
SNR. That is \emph{link adaptation}, and it is why modulation tables in modern standards read like
menus.
\end{inpractice}

\begin{history}[From 300 baud to 33.6 kb/s: the modem wars]
The analog telephone channel, 300 to 3400\,Hz wide with an SNR of perhaps 30 to 40\,dB, was the
proving ground of digital modulation for thirty years. The Bell 103 modem (1962) sent 300\,b/s by
FSK, with separate tone pairs for each direction (1070/1270\,Hz and 2025/2225\,Hz); the Bell 202
sent 1200\,b/s with 1200 and 2200\,Hz tones, and still does, in caller-ID signalling and amateur
packet radio. Phase modulation came next: the Bell 201 and 212A modems used differentially
encoded QPSK at 2400 and 1200\,b/s, and the ITU-T V.22bis modem (1984) used 16-QAM at 600 baud
for 2400\,b/s.

The leap came with coding. In 1976 Gottfried Ungerboeck, at IBM Z\"urich, proposed
\term{trellis-coded modulation} (TCM): rather than coding bits and then mapping them, design the
code and the constellation together, so that the code protects the subsets of points that lie
close together. His 1982 paper ``Channel coding with multilevel/phase signals'' showed 3 to
4\,dB of coding gain with no bandwidth expansion, and industry adopted it almost at once. V.32
(1984) sent 9600\,b/s at 2400 baud with a 32-point cross constellation and an 8-state trellis
code; V.32bis (1991) reached 14.4\,kb/s with a 128-point cross; V.34 (1994) reached 28.8 and later
33.6\,kb/s with symbol rates up to 3429 baud, four-dimensional trellis codes, precoding against
intersymbol interference, and \emph{shell mapping}, an early commercial form of the constellation
shaping of Section~\ref{sec:ch09:shaping}. V.34 operated within a few decibels of the capacity of
the telephone channel. It was the last and best analog modem; V.90 (1998) reached 56\,kb/s only by
exploiting the digital PCM network of Chapter~\ref{ch:pcm}.
\end{history}

\section{The bandwidth--power trade-off}
\label{sec:ch09:tradeoff}

\subsection{Spectral efficiency}

A linear modulation with symbol rate $R_s$ and root-raised-cosine pulses of roll-off $\beta$
occupies a passband bandwidth $W=R_s(1+\beta)$ (Chapter~\ref{ch:baseband}). Its
\term{spectral efficiency} is
\begin{equation}
\eta=\frac{R_b}{W}=\frac{\log_2M}{1+\beta}\ \ \text{b/s/Hz},
\label{eq:ch09:eta}
\end{equation}
which for ideal Nyquist pulses ($\beta=0$) is simply $\log_2M$: one bit per hertz for BPSK, two
for QPSK, six for 64-QAM. The designer's other currency is $E_b/N_0$, the energy spent per bit
relative to the noise, which through~\eqref{eq:ch03:ebno} is the received power divided by the
bit rate and $N_0$. Every modulation scheme is a point in the plane of these two currencies.

\subsection{The Shannon limit}

Shannon's capacity formula for a channel of bandwidth $W$ with white Gaussian noise
(Chapter~\ref{ch:infotheory}) is $C=W\log_2(1+P/N_0W)$. Reliable communication at rate $R_b$ is
possible if and only if $R_b<C$. Writing the received power as $P=E_bR_b$ and dividing by $W$,
\begin{equation}
\eta<\log_2\!\left(1+\eta\,\frac{E_b}{N_0}\right)
\qquad\Longleftrightarrow\qquad
\frac{E_b}{N_0}>\frac{2^\eta-1}{\eta}.
\label{eq:ch09:shannon}
\end{equation}
This is the boundary of the achievable region in the \term{bandwidth--power plane}. At
$\eta=1$ it demands $E_b/N_0>0$\,dB; at $\eta=2$, 1.76\,dB; at $\eta=6$, 10.2\,dB; at $\eta=10$,
20.1\,dB. As $\eta\to0$---unlimited bandwidth---the bound falls to
\begin{equation}
\frac{E_b}{N_0}>\ln2=0.693=-1.59\ \text{dB},
\label{eq:ch09:ln2}
\end{equation}
the \term{Shannon limit}, the minimum energy per bit with which any system, however complex and
however wasteful of bandwidth, can communicate reliably through white Gaussian noise.

\mdcfig{ch09_efficiency_plane}{The bandwidth--power plane. Each marker is an uncoded modulation
at the $E_b/N_0$ it needs for $P_b=10^{-5}$ and its spectral efficiency with ideal Nyquist
pulses; numbers are the constellation size $M$. PSK and QAM climb into the bandwidth-limited
region at a cost in power; orthogonal FSK descends towards the $-1.59$\,dB limit at a cost in
bandwidth. The 7 to 8\,dB between the uncoded points and Shannon's curve is the territory that
error-correcting codes reclaim.}

Figure~\ref{fig:ch09_efficiency_plane} places the classical modulations on this plane. Two
regimes stand out. Above $\eta=1$ is the \term{bandwidth-limited} region, where spectrum is
scarce and power relatively plentiful: cellular, Wi-Fi, cable, microwave backhaul. Designs there
climb the QAM ladder, paying about 3\,dB of $E_b/N_0$ per extra bit per hertz. Below $\eta=1$ is
the \term{power-limited} region, where every decibel of received power is precious and bandwidth
is available: deep-space probes, GPS, low-power sensors, battery-powered tags. Designs there trade
bandwidth for power, using orthogonal signalling, spreading or very-low-rate codes.

The gap between the uncoded points and the curve is striking. At $P_b=10^{-5}$ uncoded QPSK sits
7.8\,dB from the limit for $\eta=2$, and uncoded 64-QAM 7.6\,dB from it for $\eta=6$. For large QAM
constellations the gap is nearly independent of size, which leads to a handy engineering rule:
the rate of uncoded QAM is approximately
\begin{equation}
R_b\approx W\log_2\!\left(1+\frac{\mathrm{SNR}}{\Gamma}\right),
\label{eq:ch09:gap}
\end{equation}
where the \term{SNR gap} $\Gamma$ is roughly 9\,dB at symbol error rates near $10^{-6}$. The
gap approximation, due to Forney and to the DSL designers of the 1990s, is how DSL and cable
modems allocate bits across frequencies (Chapter~\ref{ch:wireline}). A code with $g$\,dB of
coding gain reduces the gap to $\Gamma-g$; modern LDPC and polar codes
(Chapter~\ref{ch:moderncodes}) cut it to 1 to 2\,dB, and the remaining 1.53\,dB of the uncoded
constellation's shape is the subject of Section~\ref{sec:ch09:shaping}.

\begin{worked}[Choosing a modulation for a satellite transponder]
A 36\,MHz satellite transponder delivers $C/N_0=90$\,dB-Hz to a ground terminal. The operator
uses RRC pulses with $\beta=0.2$, so the symbol rate is $R_s=36/1.2=30$\,Mbaud. The available
symbol SNR is
\[
\frac{E_s}{N_0}=\frac{C}{N_0}-10\log_{10}R_s=90-74.8=15.2\ \text{dB}.
\]
Uncoded, for $P_b=10^{-5}$: QPSK needs $E_s/N_0=9.6+3.0=12.6$\,dB (fits, 2.6\,dB margin);
8-PSK needs $13.0+4.8=17.8$\,dB (does not fit); 16-QAM needs $13.4+6.0=19.5$\,dB (does not fit).
Uncoded, the link carries QPSK at 60\,Mb/s. With DVB-S2 coding the picture changes completely.
On an ideal linear channel the standard's 16-APSK rate-3/4 mode needs an $E_s/N_0$ of about
10.2\,dB and 32-APSK rate-3/4 about 12.7\,dB (ETSI EN~302~307-1), so the same transponder could
carry roughly 90 or 110\,Mb/s. In practice the transponder's amplifier, driven near saturation,
eats part of that margin, and operators often settle one step lower. Chapter~\ref{ch:satellite}
develops such budgets properly, including the nonlinearity.
\end{worked}

\section{Orthogonal signalling and FSK}
\label{sec:ch09:orthogonal}

\subsection{\texorpdfstring{$M$}{M}-ary orthogonal signals}

The bottom-left of the efficiency plane belongs to \term{orthogonal signalling}: $M$ signals of
equal energy, each along its own axis, $\vect s_i=\sqrt{E_s}\,\vect e_i$. Every pair is separated
by $d=\sqrt{2E_s}$, so every other signal is a nearest neighbour, and the union bound gives
\begin{equation}
P_s\le(M-1)\,\Q\!\left(\sqrt{E_s/N_0}\right).
\end{equation}
The exact result follows from the correlation receiver. With signal 1 sent, the $M$
correlator outputs (normalised by $\sqrt{N_0/2}$) are independent Gaussians of unit variance,
one with mean $\sqrt{2E_s/N_0}$ and $M-1$ with mean zero, and a correct decision needs the first
to be the largest:
\begin{equation}
P_s=1-\int_{-\infty}^{\infty}\frac{e^{-u^2/2}}{\sqrt{2\pi}}\,
\Phi\!\left(u+\sqrt{2E_s/N_0}\right)^{M-1}du ,
\label{eq:ch09:orth}
\end{equation}
where $\Phi=1-\Q$ is the Gaussian distribution function. All errors are equally likely, and for
any bit position exactly $M/2$ of the $M-1$ wrong symbols carry the wrong bit, so
$P_b=\frac{M/2}{M-1}P_s$.

Now the remarkable part. Each orthogonal signal carries $k=\log_2M$ bits with energy
$E_s=kE_b$, but the distance between signals \emph{grows} with $k$: $d^2=2kE_b$. Adding bits to an
orthogonal set adds dimensions and pushes every signal further from every other, whereas adding
bits to QAM crowds more points into the same two dimensions. Figure~\ref{fig:ch09_orthogonal}
shows the result: the $E_b/N_0$ needed for $P_b=10^{-5}$ falls from 12.6\,dB for $M=2$ to 7.4\,dB
for $M=16$ and 4.7\,dB for $M=1024$. How far can it go? Bounding~\eqref{eq:ch09:orth} with
$\Q(x)\le\frac12e^{-x^2/2}$ gives
$P_s<M\,e^{-E_s/2N_0}=e^{-k(E_b/N_0-2\ln2)/2}$, which vanishes as $k\to\infty$ whenever
$E_b/N_0>2\ln2$ (1.42\,dB). A more careful bound, splitting the integral in two, sharpens the
threshold to exactly $\ln2$: \emph{orthogonal signalling achieves the Shannon limit
of~\eqref{eq:ch09:ln2} as $M\to\infty$}. It was the first explicit scheme known to do so. The
price is bandwidth, which grows as $M/\log_2M$, and complexity, which grows as $M$ correlators.

\mdcfig[0.85]{ch09_orthogonal}{Bit error probability of coherent $M$-ary orthogonal signalling,
from~\eqref{eq:ch09:orth}. Each doubling of $M$ lowers the required $E_b/N_0$, and as
$M\to\infty$ the curves approach a vertical wall at the Shannon limit, $-1.59$\,dB. Dotted:
noncoherent detection for $M=2$ and $M=16$.}

Two relatives are worth knowing. A \term{biorthogonal} set adds the negative of each orthogonal
signal, giving $2M$ signals in $M$ dimensions: half the bandwidth of an orthogonal set of the same
size, with almost the same power efficiency (the new antipodal pairs are further apart than the
orthogonal ones). A \term{simplex} set subtracts the centroid from an orthogonal set, saving a
fraction $1/M$ of the energy; for $M=2$ the simplex is antipodal signalling, which is why BPSK is
3\,dB better than binary orthogonal FSK. In practice orthogonal sets appear as Walsh--Hadamard
codewords (the IS-95 CDMA uplink used 64-ary orthogonal Walsh modulation with noncoherent
detection, Chapter~\ref{ch:spreadspectrum}), as pulse-position modulation in optical and
deep-space links, and as the chirps of LoRa (Chapter~\ref{ch:wifi}), which are orthogonal
cyclic shifts of a single sweep.

\subsection{Frequency-shift keying}

The most familiar orthogonal set is $M$ tones,
$s_i(t)=\sqrt{2E_s/T}\cos(2\pi f_it+\theta_i)$ with $f_i=f_0+i\,\Delta f$. How close can the tones
be and remain orthogonal? For two tones over one symbol, with $f_c\gg1/T$,
\begin{equation}
\int_0^T\cos(2\pi f_1t)\cos(2\pi f_2t+\theta)\,dt\approx
\frac{T}{2}\,\frac{\sin(2\pi\Delta fT-\theta)+\sin\theta}{2\pi\Delta fT}.
\end{equation}
If the receiver knows the phases (coherent detection, $\theta=0$), the integral vanishes when
$2\Delta fT$ is an integer: the minimum spacing is $\Delta f=1/(2T)$. If the phases are unknown
and arbitrary (noncoherent detection), it must vanish for every $\theta$, which requires
$\Delta fT$ to be an integer: the minimum spacing doubles to $\Delta f=1/T$. The occupied
bandwidth is roughly $M\Delta f$, so coherent $M$-FSK achieves
$\eta\approx2\log_2M/M$ and noncoherent $\eta\approx\log_2M/M$: one bit per hertz at best for
coherent binary or 4-ary FSK, falling rapidly for larger $M$. This is the price of the power
efficiency in Figure~\ref{fig:ch09_orthogonal}.

Binary FSK with spacing $1/(2T)$ and continuous phase has a special name---minimum-shift keying---and
a special place in Section~\ref{sec:ch09:constenv}.

\begin{history}[Teleprinters and the triumph of FSK]
The first radioteletype circuits of the 1920s and 1930s keyed a carrier on and off, like Morse
telegraphy. On fading high-frequency paths this was a nightmare: a receiver deciding ``mark'' or
``space'' by comparing the signal with a threshold had to track a signal that faded by 20 or
30\,dB within seconds. Frequency-shift keying solved it by sending a tone for mark and a different
tone for space and letting the receiver decide which tone was \emph{stronger}---a comparison that
is automatically immune to the channel's gain. That is noncoherent binary FSK, exactly the
receiver of Section~\ref{sec:ch09:noncoherent}. Commercial and military teleprinter links
standardised on shifts of 850\,Hz; radio amateurs later settled on 170\,Hz at 45.45 baud (60 words
per minute in the five-bit Baudot--Murray code). Weather broadcasts, maritime safety messages and
amateur RTTY still use these parameters, and the same idea, dressed in digital signal processing,
survives in the Bell 103 and 202 modems, in the FSK of pagers, and in the Gaussian-filtered FSK of
Bluetooth and DECT.
\end{history}

\section{Noncoherent and differential detection}
\label{sec:ch09:noncoherent}

\subsection{When the phase is unknown}

Everything so far assumed that the receiver knows the carrier phase exactly---that its local
oscillator is locked to the received carrier. Chapter~\ref{ch:sync} shows how to achieve that
and at what cost: a phase-locked loop, a training sequence or pilots, and a delay before the
first bit can be trusted. Some receivers cannot afford that cost. A burst lasting a few
milliseconds may be over before a loop settles; a channel whose phase spins with Doppler may
move faster than any loop can track; a tiny sensor may not have the power for a synthesiser
precise enough. For them, detection must work with the phase unknown.

Model the unknown phase as a random variable $\theta$ uniform on $[0,2\pi)$:
$\tilde r=\tilde s_ie^{j\theta}+\tilde w$ in complex baseband. Averaging the likelihood over
$\theta$ gives
\begin{equation}
f(\tilde r\mid i)\propto e^{-E_i/N_0}\,I_0\!\left(\frac{2\,|\langle\tilde r,\tilde s_i\rangle|}{N_0}\right),
\end{equation}
where $I_0$ is the modified Bessel function. For equal-energy signals, because $I_0$ is monotonic,
the optimal \term{noncoherent detector} picks the signal with the largest \emph{magnitude} of
correlation, $|\langle\tilde r,\tilde s_i\rangle|$, discarding the phase of the correlation along
with the unknown phase of the channel. In hardware this is a pair of correlators (I and Q) per
signal followed by an envelope $\sqrt{I^2+Q^2}$ or a square-law $I^2+Q^2$; in a digital receiver
it is the magnitude of a complex correlation, or of an FFT bin. The phase uncertainty makes one
consequence unavoidable: the detector cannot distinguish $\tilde s$ from $-\tilde s$, so
information cannot be carried in the absolute phase. Noncoherent detection works with FSK (and
orthogonal signals generally), with on-off keying, and with the phase \emph{differences} of
differential PSK.

\subsection{Noncoherent FSK}

With noncoherent binary FSK, the envelope of the correlator matched to the transmitted tone has
a Rician distribution and the envelope of the other has a Rayleigh distribution (the two curves
of Figure~\ref{fig:ch03_envelopes}). The probability that the Rayleigh one wins has a pleasingly
simple closed form,
\begin{equation}
P_b=\tfrac12\,e^{-E_b/2N_0},
\label{eq:ch09:ncfsk}
\end{equation}
and for $M$-ary orthogonal signals
\begin{equation}
P_s=\sum_{n=1}^{M-1}(-1)^{n+1}\binom{M-1}{n}\frac{1}{n+1}\exp\!\left(-\frac{n}{n+1}\,\frac{E_s}{N_0}\right).
\end{equation}
The loss relative to coherent detection is modest: for binary FSK at $10^{-5}$, 13.4\,dB against
12.6\,dB, and the loss shrinks further as $M$ grows (Figure~\ref{fig:ch09_orthogonal}). Losing the
phase costs well under a decibel---a bargain for a receiver that needs no carrier recovery at all.

\subsection{Differential PSK}

PSK needs a phase reference; \term{differential PSK} (DPSK) uses the previous symbol as one. The
transmitter encodes the data in phase \emph{changes}: $a_n=a_{n-1}\,e^{j\Delta\phi_n}$, where
$\Delta\phi_n$ is chosen from the data. The receiver forms
\begin{equation}
z_n=\tilde r_n\,\tilde r_{n-1}^{*}=|a|^2e^{j\Delta\phi_n}+\text{noise terms},
\label{eq:ch09:dpsk}
\end{equation}
in which any unknown but slowly varying channel phase cancels, and decides on the phase of $z_n$.
For binary DPSK, deciding on the sign of $\re z_n$ gives
\begin{equation}
P_b^{\text{DBPSK}}=\tfrac12\,e^{-E_b/N_0}.
\end{equation}
At $10^{-5}$ this needs 10.3\,dB against 9.6\,dB for coherent BPSK: a penalty of about 0.7\,dB,
which falls towards zero at high SNR. The penalty exists because the reference $\tilde r_{n-1}$ is
as noisy as the signal $\tilde r_n$; a detector using the previous two or more symbols
(\emph{multiple-symbol differential detection}) recovers some of it. For four phases the penalty
is larger: DQPSK with Gray labels needs 12.0\,dB, about 2.4\,dB more than coherent QPSK, and for
large $M$ the penalty approaches 3\,dB, because the effective noise, the sum of two independent
noise terms, has doubled. Figure~\ref{fig:ch09_noncoherent} compares the family.

\mdcfig[0.85]{ch09_noncoherent}{Coherent and noncoherent binary and quaternary schemes. Lines are
closed-form results; circles are simulations of differential encoding and delay-and-multiply
detection (two million bits per point). DBPSK is within a decibel of coherent BPSK; DQPSK pays
about 2.4\,dB; noncoherent FSK pays another 3\,dB for its orthogonal rather than antipodal
geometry.}

\textbf{Differential encoding}\index{differential encoding} is also used with \emph{coherent} detection, for a different reason.
A carrier recovery loop for $M$-PSK or QAM has an $M$-fold (or fourfold) phase ambiguity
(Chapter~\ref{ch:sync}): it locks with equal happiness at $0^\circ$, $90^\circ$, $180^\circ$ or
$270^\circ$. Encoding the data differentially makes the ambiguity harmless, at the cost that each
detection error corrupts two consecutive differences, roughly doubling the bit error rate---a
fraction of a decibel. Satellite and cable modems of the QPSK era almost all did this; modern
systems usually resolve the ambiguity with known pilot or synchronisation symbols instead.

\begin{inpractice}[Where differential detection lives]
Differential schemes survive wherever the channel phase is untrustworthy or receivers must be
very simple. The original 802.11 standard (1997) sent 1 and 2\,Mb/s as DBPSK and DQPSK on a
spread-spectrum carrier, and 802.11b kept those rates as its fallbacks. The European Digital
Audio Broadcasting system (DAB, ETSI EN~300~401) applies $\pi/4$-DQPSK to each OFDM subcarrier,
comparing each subcarrier with itself one symbol earlier---ideal for mobile reception, where the
channel phase changes from symbol to symbol but slowly relative to one OFDM symbol. Bluetooth's
Enhanced Data Rate modes use $\pi/4$-DQPSK (2\,Mb/s) and 8DPSK (3\,Mb/s), and the TETRA
professional mobile radio system and the second-generation American (IS-54/IS-136) and Japanese
(PDC) digital cellular standards used $\pi/4$-DQPSK. Differential detection also appears inside
coherent receivers as a quick, robust estimator of frequency offset (Chapter~\ref{ch:sync}).
\end{inpractice}

\begin{pitfall}[Differential detection is not immune to frequency offset]
DPSK tolerates an unknown \emph{phase}, but a residual carrier \emph{frequency} offset $\Delta f$
rotates $z_n$ in~\eqref{eq:ch09:dpsk} by $2\pi\Delta fT$ every symbol, biasing every decision in
the same direction. For DQPSK, whose decision regions are $\pm45^\circ$ wide, an offset of just 5\%
of the symbol rate rotates $z_n$ by $18^\circ$ and costs several decibels. Budget for it: differential
receivers still need a coarse frequency correction, and their error-rate curves assume it has been
done.
\end{pitfall}

\section{Envelope, amplifiers and continuous-phase modulation}
\label{sec:ch09:constenv}

\subsection{Why the envelope matters}

Error-rate curves compare modulations at equal \emph{average} power. Power amplifiers do not
care about average power; they care about \emph{peak} power. An amplifier driven beyond its
linear range compresses the peaks, distorting the constellation and, far worse, spreading the
signal's spectrum into neighbouring channels (the spectral regrowth of
Figure~\ref{fig:ch07_pa_regrowth}). To avoid it, a linear signal must be amplified with its
average power \emph{backed off} below the amplifier's saturation by roughly its
\term{peak-to-average power ratio} (PAPR). Back-off is expensive, because amplifier efficiency
falls with it: an ideal class-A stage converts at most 50\% of its DC power to RF at saturation
but only 12.5\% at 6\,dB back-off; an ideal class-B stage falls from 78.5\% to 39\%. In a handset,
where the PA can be the largest single consumer of battery, a few decibels of PAPR are hours of
talk time; in a satellite, they are kilograms of solar array.

Two things set a signal's PAPR: the constellation (16-QAM's corner points are 2.55\,dB above its
average) and, more importantly, the \emph{transitions} between symbols after pulse shaping. When
consecutive QPSK symbols are diametrically opposite, the filtered trajectory passes through or
near the origin, and the envelope collapses to zero before rising again
(Figure~\ref{fig:ch09_trajectories}). The resulting envelope swings produce peaks well above the
constellation's own and, through a nonlinear amplifier, strong regrowth. Figure~\ref{fig:ch09_papr_ccdf}
shows the statistics as a complementary distribution: with RRC pulses of roll-off 0.25, the power
of QPSK exceeds its average by 4.6\,dB for $10^{-4}$ of the time, 16-QAM by 6.3\,dB and 64-QAM by
6.6\,dB; an OFDM signal (Chapter~\ref{ch:ofdm}) needs about 9.6\,dB. Constant-envelope GMSK needs
none.

\mdcfig{ch09_trajectories}{I/Q trajectories of RRC-filtered ($\beta=0.35$) QPSK, offset QPSK and
$\pi/4$-QPSK, and of GMSK. Red dots mark the transmitted waveform at the symbol instants (they do
not coincide with the constellation points because the RRC pulse alone is not a Nyquist pulse;
the receiver's matched filter completes it). QPSK passes through the origin; OQPSK and
$\pi/4$-QPSK avoid the centre (dashed circle); GMSK stays on the unit circle. The minimum envelope
relative to its RMS value is printed beneath each panel.}

\mdcfig[0.85]{ch09_papr_ccdf}{Complementary cumulative distribution of instantaneous power for
several modulations with RRC pulses ($\beta=0.25$) and, for reference, an OFDM signal. The
horizontal distance at a chosen probability, commonly $10^{-4}$, is the back-off a linear
amplifier needs. GMSK is a vertical line at 0\,dB.}

\subsection{Offset QPSK and \texorpdfstring{$\pi/4$}{pi/4}-QPSK}

Two small modifications remove QPSK's zero crossings without changing its error rate in AWGN.
\textbf{Offset QPSK}\index{offset QPSK} (OQPSK, also called staggered QPSK) delays the Q stream by half a symbol,
$T_b$. I and Q never change at the same instant, so the phase jumps by at most $90^\circ$ at a time
and the trajectory never passes through the origin; the PAPR falls by about 0.6\,dB and the
envelope's minimum rises from zero to about half its RMS value. Because I and Q are still
independent BPSK streams, the matched-filter receiver and the error rate are exactly those of
QPSK. OQPSK carried the IS-95 CDMA uplink, where handset PA efficiency was paramount, and IEEE
802.15.4 (the physical layer of Zigbee and Thread) uses OQPSK with half-sine pulses at the
2.4\,GHz band---a combination we are about to recognise as MSK.

\term{$\pi/4$-QPSK} uses two QPSK constellations rotated by $45^\circ$ from each other and
alternates between them, so that every symbol changes the phase by $\pm45^\circ$ or $\pm135^\circ$.
No transition passes through the origin, and because the phase always changes, a timing
recovery circuit always has something to lock to. Encoded differentially, as
\term{$\pi/4$-DQPSK}, it can be detected either coherently or with the delay-and-multiply
receiver of Section~\ref{sec:ch09:noncoherent}. That flexibility made it the modulation of the
second-generation digital cellular systems of North America (IS-54/IS-136, 48.6\,kb/s in 30\,kHz
channels) and Japan (PDC), of the TETRA professional radio system, of DAB, and of Bluetooth EDR. The
same idea generalises. EDGE rotates 8-PSK by $3\pi/8$ per symbol to avoid the origin, and 5G NR
defines $\pi/2$-BPSK---BPSK rotated by $90^\circ$ each symbol---for its lowest-PAPR uplink mode
(3GPP TS 38.211), where it is combined with DFT-spread OFDM for cell-edge handsets.

\subsection{Continuous-phase modulation}

Removing the phase \emph{jumps} altogether leads to \term{continuous-phase modulation} (CPM),
whose complex envelope has constant amplitude and a phase that is a smooth function of the data:
\begin{equation}
\tilde s(t)=\sqrt{\frac{2E_s}{T}}\,\exp\!\Big(j\,2\pi h\sum_n a_n\,q(t-nT)\Big),
\qquad q(t)=\int_{-\infty}^tg(\tau)\,d\tau .
\label{eq:ch09:cpm}
\end{equation}
Here $a_n\in\{\pm1,\pm3,\dots\}$ are the data symbols, $h$ is the \term{modulation index}, and
$g(t)$ is the \emph{frequency pulse}, normalised so that $q(\infty)=\tfrac12$. Each symbol
therefore changes the phase by $\pi ha_n$ in total. When $g(t)$ lasts one symbol the scheme is
\emph{full response} (continuous-phase FSK); when it lasts $L>1$ symbols it is \emph{partial
response}, and each symbol's phase change is spread over several symbol periods, which narrows the
spectrum. Because the phase remembers every past symbol, a CPM signal is a finite-state
machine---its states are the accumulated phase modulo $2\pi$ and the last $L-1$ symbols---and the
optimal receiver is a Viterbi detector on that trellis (Chapter~\ref{ch:classiccodes}). Anderson,
Aulin and Sundberg's 1986 book developed this family thoroughly; its practical members are few but
enormously important.

\subsection{MSK}

\textbf{Minimum-shift keying}\index{minimum-shift keying} (MSK) is binary CPFSK with $h=\tfrac12$: a rectangular frequency pulse,
a frequency deviation of $\pm R_b/4$ (tone spacing $1/(2T_b)$, the minimum for coherent
orthogonality, hence the name), and a phase that ramps linearly by exactly $\pm90^\circ$ in every
bit (Figure~\ref{fig:ch09_cpm_phase}). It has a second, equally useful description. Working out the
I and Q components of~\eqref{eq:ch09:cpm} shows that MSK is \emph{identical} to offset QPSK in
which each bit is shaped by a half-sinusoid of duration $2T_b$:
\begin{equation}
\tilde s(t)=\sum_n a_{I,n}\,p(t-2nT_b)+j\sum_n a_{Q,n}\,p(t-(2n+1)T_b),\qquad
p(t)=\sin\frac{\pi t}{2T_b},\ 0\le t<2T_b ,
\label{eq:ch09:msk}
\end{equation}
with the bits $a_{I,n},a_{Q,n}$ related to the FSK data by a simple differential precoding. The
first view explains the constant envelope; the second explains the error rate. A linear matched-filter
receiver sees two antipodal BPSK streams, so coherent MSK achieves
$P_b=\Q(\sqrt{2E_b/N_0})$---the power efficiency of BPSK with the constant envelope of FSK. (S.
Pasupathy's 1979 tutorial, cited below, remains the clearest account of the equivalence.)

\mdcfig{ch09_cpm_phase}{Left: the phase of MSK for the bits shown ramps by $\pm90^\circ$ per bit
along the edges of the phase tree (grey); Gaussian filtering in GMSK rounds the corners and, at
$BT=0.3$, no longer reaches the full $\pm90^\circ$ within one bit. Right: the frequency pulses
$g(t)$: MSK's rectangle and GMSK's Gaussian-smoothed rectangles, which last two to three bit periods.}

The MSK spectrum falls off as $f^{-4}$, against $f^{-2}$ for QPSK with rectangular pulses, because its
phase (and hence its waveform) is continuous. Its main lobe is wider---its first null is at
$0.75R_b$ from the carrier against $0.5R_b$ for QPSK---but its 99\% power bandwidth is about
$1.2R_b$, against roughly $8R_b$ for rectangular QPSK (Figure~\ref{fig:ch09_msk_spectra}).

\subsection{GMSK}

For cellular radio even $f^{-4}$ is not enough: adjacent-channel emissions must be 60 to 70\,dB below
the carrier. \textbf{Gaussian minimum-shift keying}\index{Gaussian minimum-shift keying} (GMSK), proposed by Kazuaki Murota and Kenkichi
Hirade of NTT in 1981, passes the MSK frequency pulses through a Gaussian low-pass filter
before the frequency modulator. The filter's impulse response is
\begin{equation}
h_G(t)=\sqrt{\frac{2\pi}{\ln2}}\,B\,\exp\!\left(-\frac{2\pi^2B^2t^2}{\ln2}\right),
\end{equation}
where $B$ is its 3-dB bandwidth, and the resulting frequency pulse is
\begin{equation}
g(t)=\frac{1}{2T_b}\left[\Q\!\left(\frac{2\pi B(t-T_b/2)}{\sqrt{\ln2}}\right)-
\Q\!\left(\frac{2\pi B(t+T_b/2)}{\sqrt{\ln2}}\right)\right].
\label{eq:ch09:gmsk}
\end{equation}
The single design parameter is the normalised bandwidth $BT_b$. At $BT_b=\infty$ GMSK is MSK; as
$BT_b$ falls, the pulse spreads over more bits, the phase trajectory is smoothed, and the spectrum
narrows---at the price of controlled intersymbol interference, because each bit's phase change now
overlaps its neighbours'. At $BT_b=0.3$ the 99\% bandwidth is about $0.91R_b$ and the out-of-band
spectrum falls much faster than MSK's (Figure~\ref{fig:ch09_msk_spectra}); the ISI costs a fraction
of a decibel in a receiver that accounts for it. GSM transmits 270.833\,kb/s of GMSK with
$BT_b=0.3$ in 200\,kHz channels, about 1.35\,b/s/Hz. Because GMSK is so close to linear, its
receivers usually treat it as a linear modulation using the first term of Laurent's 1986
decomposition of CPM into a sum of amplitude-modulated pulses, and absorb the residual ISI into the
Viterbi equaliser that multipath requires anyway (Chapter~\ref{ch:equalization}).

\mdcfig[0.85]{ch09_msk_spectra}{Power spectral densities (normalised to unit total power) against
frequency offset in units of the bit rate. Rectangular QPSK has nulls every $R_b/2$ and slowly
decaying sidelobes; MSK's continuous phase gives $f^{-4}$ sidelobes; Gaussian filtering with
$BT=0.5$ and $0.3$ suppresses them further. RRC-filtered QPSK is the most compact of all but has a
non-constant envelope.}

\begin{history}[How GSM chose GMSK]
In the mid-1980s the European Conference of Postal and Telecommunications Administrations had to
choose a radio interface for a pan-European digital cellular system, and several consortia
proposed candidates ranging from wideband designs to narrowband TDMA. After comparative
field trials in Paris in 1986, the Groupe Sp\'ecial Mobile agreed in 1987 on narrowband TDMA with
eight time slots per 200\,kHz carrier and GMSK modulation. The reasons were engineering ones.
GMSK's constant envelope let handsets use efficient, saturated power amplifiers at a time when
linear amplifiers were costly and power-hungry; its tight spectrum allowed 200\,kHz channel
spacing; and its noise immunity, close to that of BPSK, suited a cellular system limited by
co-channel interference. The price was spectral efficiency, which EDGE later recovered by
switching the same carriers to 8-PSK. By the time GSM was retired in many countries in the 2020s,
several billion GMSK handsets had been built.
\end{history}

\begin{inpractice}[The GFSK family: Bluetooth, BLE and DECT]
When the modulation index departs from $\tfrac12$ the scheme is Gaussian FSK (GFSK), which trades
some power efficiency for tolerance of cheap, imprecise oscillators and for simple
frequency-discriminator receivers. Classic Bluetooth (Basic Rate) sends 1\,Mb/s of GFSK with
$BT=0.5$ and a modulation index between 0.28 and 0.35. Bluetooth Low Energy uses GFSK with
$BT=0.5$ and index 0.45 to 0.55 at 1 or 2\,Mb/s. DECT cordless telephones use GFSK with $BT=0.5$
and $h=0.5$ at 1.152\,Mb/s. The maritime Automatic Identification System uses GMSK at 9.6\,kb/s.
All share the virtues that made GSM choose GMSK: a saturated transmitter, a spectrum that dies
away quickly, and a receiver that can be a limiter and a discriminator costing a few square
millimetres of silicon.
\end{inpractice}

\begin{keyidea}[Constant envelope versus spectral efficiency]
Constant-envelope modulations (MSK, GMSK, GFSK) let the power amplifier run in saturation at high
efficiency but are limited to about 1 to 1.5\,b/s/Hz. Linear modulations (QAM with RRC pulses)
reach many bits per hertz but need a linear, backed-off amplifier. Offset and rotated schemes sit
between. Modern handsets resolve the dilemma with digital predistortion and envelope tracking
(Chapter~\ref{ch:sdr}), which make linear amplification efficient enough for QAM and OFDM, but the
trade-off still decides the modulation of every battery-powered sensor and satellite.
\end{keyidea}

\section{Soft decisions: bit log-likelihood ratios}
\label{sec:ch09:soft}

\subsection{From hard to soft}

A detector that outputs only the nearest symbol throws away information: a sample that lands
exactly on a constellation point and one that lands on a decision boundary yield the same hard
decision, but the second is a coin toss. Modern error-correcting decoders
(Chapters~\ref{ch:classiccodes} and~\ref{ch:moderncodes}) can use that reliability information,
and it is worth about 2\,dB to them. The standard currency is the bit \term{log-likelihood ratio}
(LLR). For the $k$th bit of the label of the received symbol,
\begin{equation}
L(b_k)=\ln\frac{P(b_k=0\mid y)}{P(b_k=1\mid y)}
=\ln\frac{\sum_{x\in\mathcal X_k^0}p(x)\,e^{-|y-x|^2/N_0}}{\sum_{x\in\mathcal X_k^1}p(x)\,e^{-|y-x|^2/N_0}},
\label{eq:ch09:llr}
\end{equation}
where $\mathcal X_k^b$ is the set of constellation points whose $k$th bit equals $b$ and $p(x)$ is
the prior of each point (uniform in most systems). The sign of $L$ is the hard decision and its
magnitude the confidence: $|L|=0$ means no idea, $|L|=5$ means the decision is wrong with
probability $1/(1+e^5)\approx0.7\%$.

Evaluating~\eqref{eq:ch09:llr} for 1024-QAM needs 1024 exponentials per symbol. The
\term{max-log approximation} keeps only the largest term in each sum:
\begin{equation}
L(b_k)\approx\frac{1}{N_0}\Big[\min_{x\in\mathcal X_k^1}|y-x|^2-\min_{x\in\mathcal X_k^0}|y-x|^2\Big],
\label{eq:ch09:maxlog}
\end{equation}
the difference between the squared distances to the nearest point with the bit equal to 1 and
the nearest with it equal to 0. For Gray-labelled square QAM the minimisations separate by axis
and reduce to piecewise-linear functions of the I or Q sample. For 16-QAM with levels
$\{-3,-1,+1,+3\}$ on each axis, labels $00,01,11,10$ from left to right, and noise variance $N_0/2$
per dimension, the in-phase bits are
\begin{equation}
L(b_1)\approx
\begin{cases}-4y_I/N_0, & |y_I|\le2\\ -8(y_I-\operatorname{sgn}y_I)/N_0,& |y_I|>2\end{cases},
\qquad
L(b_2)\approx\frac{4(|y_I|-2)}{N_0},
\label{eq:ch09:llr16}
\end{equation}
which a receiver evaluates with a few adds and shifts (Figure~\ref{fig:ch09_llr}).

\mdcfig{ch09_llr}{Left: exact (solid) and max-log (dashed) LLRs of the two in-phase bits of
Gray-labelled 16-QAM at $E_s/N_0=5$\,dB, against the in-phase sample; dotted lines mark the
constellation levels. The approximation is visibly inexact only near the decision boundaries.
Right: distribution of the LLR, signed so that positive means correct, at 12\,dB. The first
(sign) bit makes half as many errors as the second: bit positions in a QAM symbol are not equal.}

\begin{worked}[Soft bits from one 16-QAM sample]
A receiver with $N_0=0.5$ (on the unnormalised $\pm1,\pm3$ grid) observes $y=0.5-2.4j$.
Using~\eqref{eq:ch09:llr16} for I and the same expressions for Q:
\begin{itemize}
\item I: $L(b_1)=-4(0.5)/0.5=-4.0$; $L(b_2)=4(0.5-2)/0.5=-12.0$.
\item Q: $L(b_3)=-8(-2.4+1)/0.5=+22.4$; $L(b_4)=4(2.4-2)/0.5=+3.2$.
\end{itemize}
The hard decisions are bits $1,1,0,0$: I level $+1$, Q level $-3$, which is indeed the nearest
point, $(1,-3)$. But the confidences differ enormously. The third bit is essentially certain
(error probability $\approx2\times10^{-10}$); the fourth, from a sample only 0.4 beyond the Q
boundary at $-2$, is wrong with probability $1/(1+e^{3.2})\approx4\%$. A hard-decision receiver
would pass all four bits to the decoder as equally trustworthy; a soft one tells the decoder
exactly where to look for errors.
\end{worked}

\subsection{Bit-interleaved coded modulation}

Soft demapping makes possible the architecture used by essentially every modern standard---Wi-Fi,
LTE and NR, DVB-S2/T2, DOCSIS 3.1: \term{bit-interleaved coded modulation} (BICM). A binary code
(LDPC, turbo or polar) encodes the data; an interleaver scatters the coded bits; a Gray mapper
groups them into symbols; the receiver computes LLRs with~\eqref{eq:ch09:llr}
or~\eqref{eq:ch09:maxlog} and de-interleaves them into a binary decoder. The interleaver makes the
bits of each symbol, which have different reliabilities (Figure~\ref{fig:ch09_llr}) and fade
together, look to the decoder like independent binary channels. BICM was introduced by Ephraim
Zehavi in 1992 for fading channels and analysed by Caire, Taricco and Biglieri in 1998, who showed
that with Gray labelling it loses only a fraction of a decibel relative to the best joint
coding-and-modulation schemes---while letting one binary code serve every constellation. It
displaced Ungerboeck's trellis-coded modulation in wireless systems for exactly that reason. The
theory of these capacities is developed in Chapter~\ref{ch:infotheory}.

\begin{pitfall}[LLRs need the right noise level]
Equation~\eqref{eq:ch09:llr} contains $N_0$. A demapper fed a wrong noise estimate produces LLRs
that are all too large (overconfident) or too small. Min-sum LDPC decoders are nearly insensitive to
a common scale factor, but belief-propagation, turbo and polar list decoders are not, and
overconfident LLRs can cost a decibel or more. With QAM there is a second trap: the demapper needs
the \emph{amplitude} reference---the channel gain $|h|$ that scales the grid---as well as the phase.
A demapper that uses an unscaled grid after a fade computes the distances to the wrong points.
Channel estimation (Chapter~\ref{ch:ofdm}) must supply both.
\end{pitfall}

\section{A first look at fading}
\label{sec:ch09:fading}

Everything so far has assumed a fixed channel gain. A mobile radio channel is not so kind:
multipath components add with phases that change as the terminal moves a fraction of a
wavelength, and the received amplitude fluctuates by tens of decibels (Chapter~\ref{ch:channels}).
The simplest model, \term{Rayleigh fading}, takes the channel gain $h$ in $y=ha+n$ to be a
circularly symmetric complex Gaussian, so $|h|$ is Rayleigh distributed and the instantaneous SNR
$\gamma=|h|^2\bar\gamma$ is exponentially distributed about its mean $\bar\gamma$. If the fading is
slow enough to estimate, the receiver detects coherently with the known $h$, and the error rate at
each instant is the AWGN value at that instant's SNR. Averaging over the exponential distribution
gives, for BPSK,
\begin{equation}
\bar P_b=\int_0^\infty\Q\!\left(\sqrt{2\gamma}\right)\frac{e^{-\gamma/\bar\gamma}}{\bar\gamma}\,d\gamma
=\frac12\left(1-\sqrt{\frac{\bar\gamma}{1+\bar\gamma}}\right)\approx\frac{1}{4\bar\gamma},
\label{eq:ch09:rayleigh}
\end{equation}
and for DBPSK $1/(2(1+\bar\gamma))$ and noncoherent FSK $1/(2+\bar\gamma)$.

\mdcfig{ch09_rayleigh}{Bit error probability in flat Rayleigh fading (solid) against AWGN (dashed).
The waterfall becomes a straight line falling one decade per 10\,dB. Circles are simulations of
BPSK with a perfectly known channel. Maximal-ratio combining of two or four independently fading
branches (dash-dot, plotted against the SNR per branch) restores a steep slope.}

The change is dramatic (Figure~\ref{fig:ch09_rayleigh}). The exponential waterfall of the AWGN
channel becomes a straight line falling only one decade per 10\,dB. BPSK needs 6.8\,dB for
$10^{-3}$ in AWGN but 24\,dB in Rayleigh fading; for $10^{-5}$ the figures are 9.6 and 44\,dB, a
34\,dB penalty. The cause is not the average SNR but the occasional deep fade: the channel is
more than 10\,dB below its mean about 10\% of the time and more than 20\,dB below about 1\% of the
time, and during those fades errors are almost certain. Bigger constellations do not escape; they
suffer the same inverse-SNR law with a larger constant.

The cure is \term{diversity}: give the receiver several independently fading copies of each bit,
so that all of them are unlikely to fade at once. With $L$ independent branches combined
optimally (maximal-ratio combining, Chapter~\ref{ch:mimo}), the error rate falls as
$\bar\gamma^{-L}$: BPSK at $10^{-3}$ needs 24\,dB per branch with one branch, 11.1\,dB with two and
4.0\,dB with four. The branches can be antennas (spatial diversity), frequencies (OFDM with coding
across subcarriers, Chapter~\ref{ch:ofdm}), times (coding with interleaving deeper than a fade) or
paths (the rake receiver of Chapter~\ref{ch:spreadspectrum}). Almost every design decision in a
mobile physical layer is, at bottom, a way of buying diversity.

\begin{keyidea}[Fading changes what matters]
In AWGN, performance is set by minimum Euclidean distance and a code's job is to increase it. In
fast Rayleigh fading with interleaving, performance is set by \emph{diversity order}---how many
independent fades must coincide to cause an error---and the relevant property of a code is the
minimum \emph{Hamming} distance between codewords in symbols, which is exactly what BICM's bit
interleaver maximises. Chapter~\ref{ch:channels} develops the models, Chapter~\ref{ch:mimo} the
antenna techniques.
\end{keyidea}

\section{Measuring modulation quality: EVM and MER}
\label{sec:ch09:evm}

\subsection{Error vector magnitude}

Error-rate curves describe an ideal transmitter and an ideal receiver separated by noise. Real
transmitters are not ideal: their oscillators wander, their I and Q paths are mismatched, their
amplifiers compress, their DACs quantize. The industry's universal yardstick for all of these at
once is the \term{error vector magnitude} (EVM). A test receiver---a vector signal analyser---
demodulates the signal, synchronises it, and compares each received symbol $y_n$ with the ideal
constellation point $s_n$ that it represents. The difference $e_n=y_n-s_n$ is the error vector,
and
\begin{equation}
\mathrm{EVM}_{\text{rms}}=\sqrt{\frac{\sum_n|y_n-s_n|^2}{\sum_n|s_n|^2}},\qquad
\mathrm{EVM}_{\text{dB}}=20\log_{10}\mathrm{EVM}_{\text{rms}} .
\label{eq:ch09:evm}
\end{equation}
EVM is quoted in percent (cellular) or decibels (Wi-Fi); 3.5\% is $-29.1$\,dB. Broadcasting and cable
engineers use the reciprocal, the \term{modulation error ratio}
$\mathrm{MER}=10\log_{10}(\sum|s_n|^2/\sum|e_n|^2)$ (ETSI TR~101~290), which equals $-\mathrm{EVM}_{\text{dB}}$
when both are normalised to the average constellation power.

If the error vectors are noise-like, EVM is simply an SNR in disguise: $\mathrm{SNR}\approx1/\mathrm{EVM}^2$.
This makes it additive with everything else in a link budget. A transmitter with EVM
$\epsilon_{\text{tx}}$ delivers to even a noiseless receiver an effective SNR of
$1/\epsilon_{\text{tx}}^2$, and a receiver with thermal SNR $\rho$ sees
\begin{equation}
\frac{1}{\mathrm{SNR}_{\text{eff}}}=\frac1\rho+\epsilon_{\text{tx}}^2+\epsilon_{\text{rx}}^2 .
\label{eq:ch09:evmsnr}
\end{equation}
A 1024-QAM link needs an effective SNR in the thirties of decibels; a transmitter with 3\% EVM
($-30.5$\,dB) caps it at 30.5\,dB however strong the signal. That is why EVM limits, not noise,
decide how large a constellation a short-range system can use.

\subsection{What impairments look like}

Each impairment leaves a characteristic signature on the constellation, and an experienced
engineer diagnoses a transmitter from its constellation diagram the way a physician reads an
X-ray (Figure~\ref{fig:ch09_impairments}):
\begin{itemize}
\item \textbf{Additive noise} makes circular clouds of equal size on every point.
\item \textbf{Phase noise} (Chapter~\ref{ch:sdr}) makes arcs, longer on the outer points because a
phase error $\delta\phi$ moves a point of amplitude $A$ by $A\,\delta\phi$. For small errors
$\mathrm{EVM}\approx\sigma_\phi$ in radians: $1^\circ$ rms costs 1.7\%, $3^\circ$ costs 5.2\%.
\item \textbf{IQ imbalance}---unequal gain or a non-$90^\circ$ phase between the I and Q paths---skews
the grid into a parallelogram, because the received signal is $\mu s+\nu s^*$ with an image term
$\nu s^*$. The error it causes is set by the image rejection ratio $|\mu/\nu|^2$: 30\,dB of image
rejection limits EVM to about $-30$\,dB.
\item \textbf{Amplifier compression} pulls the outer points inwards and, with AM/PM, twists them;
because it acts on the filtered waveform, not the symbols, it also smears every point into a cloud
and spreads the spectrum.
\item \textbf{Residual frequency offset} rotates the constellation progressively through a burst,
turning the clouds into arcs that grow with time.
\item \textbf{Carrier leakage} (LO feedthrough or DC offset) shifts the whole constellation off
centre; most standards measure and limit it separately from EVM.
\end{itemize}

\mdcfig{ch09_impairments}{The signatures of common impairments on 16-QAM, each with a small amount
of additive noise and with the raw EVM of~\eqref{eq:ch09:evm}. A measuring instrument would first
remove the average gain, phase and frequency offset---standards specify exactly which corrections
are allowed---so the IQ-imbalance and frequency-offset cases would read lower on a real analyser.}

\begin{pitfall}[Read the EVM definition before comparing numbers]
EVM figures from different sources are often not comparable. Standards normalise to average power
(3GPP, IEEE 802.11) or, in some older specifications, to peak constellation power, which makes the
same signal look better by the constellation's peak-to-average ratio (2.55\,dB for 16-QAM,
4.23\,dB for 256-QAM). They differ in what the analyser may correct first: 3GPP TS 38.104 measures
EVM after a defined equaliser and reports frequency error separately; 802.11 estimates the channel
from the preamble and tracks phase with pilots. Some quote the average over many bursts, some the
worst burst. And an EVM measured with a noisy receiver includes that receiver's own EVM
through~\eqref{eq:ch09:evmsnr}. Always state the definition, the correction and the measurement
bandwidth alongside the number.
\end{pitfall}

\subsection{EVM in the standards}

Standards specify a maximum transmitter EVM for each modulation, chosen so that the transmitter
consumes only a small part of the SNR the receiver needs. For 5G NR base stations, 3GPP TS~38.104
requires an EVM of at most 17.5\% for QPSK, 12.5\% for 16-QAM, 8\% for 64-QAM, 3.5\% for 256-QAM and
2.5\% for 1024-QAM; the handset limits in TS~38.101-1 follow the same pattern up to 256-QAM.
Converted to decibels these are $-15.1$, $-18.1$, $-21.9$, $-29.1$ and $-32.0$\,dB. Wi-Fi's limits,
which are tighter because Wi-Fi uses higher code rates with its large constellations, appear in
Table~\ref{tab:ch09:wifimcs}.

\begin{table}[htbp]
\centering
\caption{The IEEE 802.11ax (Wi-Fi 6) modulation and coding schemes, with the two 4096-QAM schemes
added by 802.11be (Wi-Fi 7). Data rates are for one spatial stream in a 20\,MHz channel (234 data
subcarriers) with the 0.8\,$\mu$s guard interval. The transmitter constellation error limits are those
of IEEE Std 802.11ax-2021 and, for MCS~12--13, IEEE 802.11be.}
\label{tab:ch09:wifimcs}
\small
\begin{tabular}{clcccr}
\toprule
MCS & Modulation & Code rate & Bits/subcarrier & Data rate (Mb/s) & Max.\ EVM (dB)\\
\midrule
0 & BPSK & 1/2 & 1 & 8.6 & $-5$\\
1 & QPSK & 1/2 & 2 & 17.2 & $-10$\\
2 & QPSK & 3/4 & 2 & 25.8 & $-13$\\
3 & 16-QAM & 1/2 & 4 & 34.4 & $-16$\\
4 & 16-QAM & 3/4 & 4 & 51.6 & $-19$\\
5 & 64-QAM & 2/3 & 6 & 68.8 & $-22$\\
6 & 64-QAM & 3/4 & 6 & 77.4 & $-25$\\
7 & 64-QAM & 5/6 & 6 & 86.0 & $-27$\\
8 & 256-QAM & 3/4 & 8 & 103.2 & $-30$\\
9 & 256-QAM & 5/6 & 8 & 114.7 & $-32$\\
10 & 1024-QAM & 3/4 & 10 & 129.0 & $-35$\\
11 & 1024-QAM & 5/6 & 10 & 143.4 & $-35$\\
\midrule
12 & 4096-QAM & 3/4 & 12 & 154.9 & $-38$\\
13 & 4096-QAM & 5/6 & 12 & 172.1 & $-38$\\
\bottomrule
\end{tabular}
\end{table}

\begin{worked}[Reading the Wi-Fi MCS table]
An 802.11ax OFDM symbol in a 20\,MHz channel lasts $12.8\,\mu$s plus a $0.8\,\mu$s guard interval,
13.6\,$\mu$s in all, and carries data on 234 subcarriers (Chapter~\ref{ch:ofdm}). At MCS~11 each
subcarrier carries a 1024-QAM symbol, 10 coded bits, of which a fraction 5/6 are information:
\[
R=\frac{234\times10\times5/6}{13.6\,\mu\text{s}}=\frac{1950\ \text{bits}}{13.6\,\mu\text{s}}=143.4\ \text{Mb/s}.
\]
Doubling the bandwidth roughly doubles the rate (a 160\,MHz channel has 1960 data subcarriers, giving
1201\,Mb/s per stream), and eight spatial streams multiply it by eight, which is where the headline
figure of 9.6\,Gb/s for Wi-Fi~6 comes from. The EVM column shows what the headline costs: a
transmitter whose EVM is $-35$\,dB has already spent the SNR budget of a perfect receiver at 35\,dB,
and the receiver's own noise figure, phase noise and quantization come on top
via~\eqref{eq:ch09:evmsnr}. In practice MCS~10--11 is reachable only within a few metres of the
access point, with a clear line of sight.
\end{worked}

\begin{worked}[An EVM budget for a 1024-QAM transmitter]
A designer must meet $-35$\,dB (1.78\%) for 1024-QAM. The contributions, estimated or measured
separately and treated as independent, are: synthesiser phase noise of $0.5^\circ$ rms
($\sigma_\phi=8.7\times10^{-3}$\,rad, $-41.2$\,dB); IQ image rejection of 45\,dB after calibration
($-45$\,dB); power amplifier nonlinearity at the planned back-off, with digital predistortion,
$-38$\,dB; DAC quantization, clock jitter and analog noise, $-40$\,dB. Error powers add:
\[
10^{-4.12}+10^{-4.5}+10^{-3.8}+10^{-4.0}=(0.76+0.32+1.58+1.00)\times10^{-4}=3.66\times10^{-4},
\]
which is $-34.4$\,dB, or 1.91\%. The design fails by 0.6\,dB, and the table of contributions shows
where to act: the PA term dominates. Backing the PA off by 1\,dB typically improves its
intermodulation-limited EVM by about 2\,dB (third-order products fall twice as fast as the
signal), to $-40$\,dB; the total becomes $3.08\times10^{-4}$, or $-35.1$\,dB. It passes, with no
margin, at the cost of about a decibel of output power. Real designs keep 2 to 3\,dB of margin for
production spread and temperature.
\end{worked}

\section{Constellation shaping}
\label{sec:ch09:shaping}

\subsection{The 1.53\,dB that uniform constellations leave behind}

Error-correcting codes close most of the 7 to 8\,dB gap between uncoded QAM and the Shannon limit,
but not all of it. Part of the gap has nothing to do with coding: it comes from the \emph{shape} of
the constellation. Shannon's capacity is achieved by Gaussian-distributed channel inputs, whereas
a square QAM constellation with equiprobable points approximates a \emph{uniform} distribution over
a square. At high SNR the cost can be computed exactly. A uniform distribution on an interval of
width $a$ has differential entropy $\log_2a$ and variance $a^2/12$; a Gaussian with the same
entropy has variance $a^2/(2\pi e)$. The uniform input therefore spends
\begin{equation}
\frac{a^2/12}{a^2/2\pi e}=\frac{\pi e}{6}=1.423\quad(1.53\ \text{dB})
\label{eq:ch09:shaping}
\end{equation}
more energy per dimension than necessary to carry the same information. This \term{shaping gain}
of 1.53\,dB is the most that any rearrangement or reweighting of the constellation can recover,
and it is independent of the code: a perfect code on a uniform square constellation stays 1.53\,dB
from capacity at high SNR. The ideal is the familiar picture of Shannon's proof: points packed
uniformly inside a \emph{sphere} of high dimension rather than a cube, whose projection onto any
one dimension is Gaussian.

There are two ways to get the gain. \textbf{Geometric shaping}\index{geometric shaping} moves the points: rings with more
points on the inside (APSK), non-uniform grids, or, as in the V.34 modem's shell mapping, choosing
multidimensional points from inside a sphere rather than a cube. ATSC~3.0 broadcasting and DVB-S2X
use such non-uniform constellations, optimised numerically for each code rate. \textbf{Probabilistic
shaping}\index{probabilistic shaping} keeps the regular grid but sends the points with unequal probabilities---inner points
often, outer points rarely---so that the input distribution approximates a sampled Gaussian. The
natural choice is the \term{Maxwell--Boltzmann distribution}
\begin{equation}
P(x)=\frac{e^{-\lambda|x|^2}}{\sum_{x'}e^{-\lambda|x'|^2}},
\label{eq:ch09:mb}
\end{equation}
which maximises entropy for a given average energy (Kschischang and Pasupathy, 1993). The single
parameter $\lambda$ trades rate against energy: $\lambda=0$ is uniform QAM, and increasing it lowers
both the entropy and the average power.

\mdcfig{ch09_shaping}{Probabilistic shaping. Left: a Maxwell--Boltzmann distribution on 64-QAM; dot
area is proportional to probability, and the entropy of 5.57\,bits is less than the 6\,bits of
uniform 64-QAM. Centre: mutual information (achievable rate) of uniform 16-, 64- and 256-QAM and of
256-QAM with the Maxwell--Boltzmann parameter optimised at each SNR. Right: the SNR gap to Shannon's
formula at a given rate. Uniform QAM approaches the 1.53\,dB shaping loss at high rates; the shaped
constellation stays within about 0.1\,dB of capacity until its 8-bit ceiling approaches.}

Figure~\ref{fig:ch09_shaping}, computed with numerical integration of the mutual information, shows
the payoff. At a rate of 6 bits per symbol uniform 256-QAM needs an SNR of 19.2\,dB, 1.2\,dB above
Shannon's 18.0\,dB; Maxwell--Boltzmann-shaped 256-QAM needs 18.1\,dB. At 4 bits per symbol the
gain is 0.8\,dB. Shaping also delivers a second, commercially just as important, benefit: because
$\lambda$ is continuous, the spectral efficiency can be tuned in arbitrarily fine steps with a single
code and constellation, rather than in jumps between modulation formats.

\subsection{Probabilistic amplitude shaping}

The practical obstacle was always the interface with coding: a binary encoder produces
equiprobable bits, which a mapper turns into equiprobable points. The \term{probabilistic amplitude
shaping} (PAS) architecture of B\"ocherer, Steiner and Schulte (2015) removed it with a neat
observation. The Maxwell--Boltzmann distribution depends only on $|x|$, so the \emph{signs} of the
I and Q coordinates can remain uniform. A \emph{distribution matcher} (for example, constant
composition distribution matching) converts uniform data bits into a sequence of amplitudes with
the desired distribution; the amplitudes' bit labels pass through a systematic binary LDPC encoder,
whose systematic output is unchanged and whose parity bits---being uniform---are used as the
signs. The receiver computes LLRs with the prior term of~\eqref{eq:ch09:llr}, decodes, and inverts
the matcher. The existing binary FEC is untouched.

\begin{inpractice}[Shaping in coherent optics]
Probabilistic shaping made the leap from theory to product faster than almost any idea in this
book. Coherent optical transceivers carry dual-polarisation QAM over fibre at 60 to 150\,Gbaud,
where every fraction of a decibel of optical SNR translates into reach---tens or hundreds of
kilometres before a regenerator is needed. The first commercial probabilistically shaped DSP
(Nokia's PSE-3) was announced in 2019, and shaped constellations have since become a standard
feature of high-end long-haul transceivers, where they deliver on the order of 1\,dB of margin and
continuous rate adaptation from roughly 200\,Gb/s to 800\,Gb/s per wavelength. Wireless standards
have been slower to follow, because shaping complicates rate matching and hybrid ARQ and raises
the peak-to-average ratio (outer points are rarer, so the average falls while the peak stays put);
it is a leading candidate for the sixth generation.
\end{inpractice}

\begin{keyidea}[Where the decibels go]
Between uncoded QAM at $10^{-5}$ and the Shannon limit lie roughly 7 to 8\,dB. Coding recovers most
of it: about 5 to 6\,dB for modern LDPC or polar codes with soft decisions. Shaping recovers up to
another 1.53\,dB. What remains---a few tenths of a decibel for long codes---is the price of finite
block length and practical decoders. Every one of those decibels is a factor of 1.26 in range in free
space, or in battery life, or in satellite antenna area.
\end{keyidea}

\section{Summary}
\label{sec:ch09:summary}

\begin{itemize}
\item Waveforms are vectors. $M$ signals in white Gaussian noise are described completely by their
coordinates in an $N\le M$ dimensional signal space, and the projections of the received signal onto
that space are a sufficient statistic. For linear modulation $N=2$: one I and one Q sample.
\item The MAP detector minimises $\|\vect r-\vect s_i\|^2-N_0\ln p_i$; with equal priors it is the
minimum-distance detector, whose decision regions are Voronoi cells. It can be built from
correlators or matched filters.
\item Error probability is governed by distances: $P(\vect s_i\to\vect s_j)=\Q(\sqrt{d_{ij}^2/2N_0})$,
and at useful SNR $P_s\approx\bar N_{\min}\Q(\sqrt{d_{\min}^2/2N_0})$. With Gray labelling
$P_b\approx P_s/\log_2M$.
\item QPSK is as power-efficient as BPSK at twice the rate. Beyond eight points QAM beats PSK; each
doubling of QAM size costs about 3\,dB of $E_s/N_0$ per extra bit, and uncoded 64-QAM needs 17.8\,dB
of $E_b/N_0$ for $10^{-5}$. APSK trades a little distance for a lower PAPR on satellite amplifiers.
\item In the bandwidth--power plane, Shannon's limit is $E_b/N_0>(2^\eta-1)/\eta$, reaching $-1.59$\,dB
as $\eta\to0$. Orthogonal signalling approaches it with unlimited bandwidth; uncoded QAM sits 7 to
8\,dB away.
\item Noncoherent and differential detection avoid carrier recovery for well under a decibel (FSK,
DBPSK) to about 2.4\,dB (DQPSK).
\item Envelope variations cost PA efficiency. OQPSK and $\pi/4$-QPSK avoid the origin; MSK and GMSK
have a constant envelope, MSK with the error rate of BPSK and GMSK ($BT=0.3$ in GSM) with a compact
spectrum.
\item Decoders want soft bits: LLRs, computed exactly or with the max-log approximation, feed the
bit-interleaved coded modulation used by every modern standard.
\item Rayleigh fading turns the exponential waterfall into a $1/\mathrm{SNR}$ law; diversity restores it.
\item EVM and MER measure transmitter quality; they add to receiver noise as an SNR limit. Shaping
recovers up to 1.53\,dB by making constellations look Gaussian.
\end{itemize}

\begin{labbox}[Signal space, detection and constellations]
\raggedright
Lab~2, \lab{moderncomms-labs/labs/lab02\_modulation.py} (notebook
\lab{lab02\_modulation.ipynb}), is the companion to this chapter. It builds the BPSK to 256-QAM
constellations with \texttt{commlib.get\_constellation} and tabulates their $d_{\min}$ and PAPR;
colours the ML decision regions of 8-PSK and 16-QAM; measures bit error rates by Monte Carlo for
QPSK, 8-PSK, 16-QAM and 64-QAM and overlays the closed forms of this chapter; offers an interactive
explorer (constellation and $E_b/N_0$ sliders) that reports BER, SER and EVM live; compares Gray and
natural labelling on 16-QAM; and plots exact and max-log LLRs for the 16-QAM bits. Its exercises
ask you to derive $d_{\min}$ for square QAM, check the 8-PSK union bound, find where 8-PSK and 16-QAM
cross, and optimise the ring ratio of a DVB-S2 style 16-APSK. Every figure in this chapter is
produced by \texttt{book/figscripts/ch09\_figs.py} using the same library; the GMSK generator, the
impairment models and the mutual-information integrator there are good starting points for the
simulation problems below. A second lab, \lab{moderncomms-labs/labs/lab21\_cpm\_evm.py},
covers MSK/GMSK generation and detection (coherent, differential and discriminator receivers),
occupied bandwidth, PAPR, and EVM measurement of simulated transmitter impairments.
\end{labbox}

\problems
\begin{enumerate}[label=\thechapter.\arabic*]
\item Four signals on $[0,3)$ are defined piecewise on the unit intervals $[0,1),[1,2),[2,3)$ by the
values $s_1=(1,1,0)$, $s_2=(1,-1,1)$, $s_3=(2,0,1)$ and $s_4=(0,2,-1)$. Apply Gram--Schmidt, find the
dimension of the signal space, the coordinates of each signal and all six pairwise distances.
Which pair is most likely to be confused in AWGN?
\item \emph{(Priors)} Binary antipodal signals $\pm1$ are received in Gaussian noise of standard
deviation $\sigma$ with $P(-1)=0.9$. Derive the MAP threshold~\eqref{eq:ch09:mapthresh}. For
$\sigma=0.5$ and $\sigma=1$, compute the error probability with the ML and MAP thresholds. Explain
why the gain vanishes at high SNR.
\item Show that the minimum distance of square $M$-QAM with unit average energy is
$\sqrt{6/(M-1)}$, and that the average number of nearest neighbours of 64-QAM is 3.5. Use the
nearest-neighbour approximation to estimate the SER of 64-QAM at $E_s/N_0=24$\,dB and compare with
the exact expression~\eqref{eq:ch09:qam}.
\item For 8-PSK, write the union bound including the first and second nearest neighbours. Evaluate
it at $E_s/N_0=10$ and 14\,dB and compare with Craig's formula~\eqref{eq:ch09:craig}. At which SNR is
the nearest-neighbour term within 10\% of the exact result?
\item Estimate the $E_b/N_0$ needed by uncoded 256-QAM for $P_b=10^{-6}$. How far is it from the
Shannon limit at $\eta=8$\,b/s/Hz? If a code with 6\,dB of coding gain at rate 7/8 is added, what
spectral efficiency and $E_b/N_0$ result?
\item Verify the values of $d_{\min}^2/E_s$ quoted for 32-cross QAM (0.200) and the $4\times8$
rectangle (0.154). Propose a labelling of the 32-cross in which as few neighbour pairs as possible
differ in more than one bit, and count them.
\item Derive the minimum tone spacings for coherent ($1/2T$) and noncoherent ($1/T$) orthogonal FSK.
What are the approximate spectral efficiencies of coherent and noncoherent 8-FSK and 64-FSK?
\item Show that MSK, defined as CPFSK with $h=\tfrac12$ and rectangular frequency pulses, can be
written in the OQPSK form~\eqref{eq:ch09:msk} with half-sine pulses. Find the relationship between
the FSK data bits and the I/Q bits.
\item \emph{(DPSK)} Show that binary DPSK over two symbol periods is equivalent to binary orthogonal
signalling with energy $2E_b$, detected noncoherently, and hence derive $P_b=\tfrac12e^{-E_b/N_0}$ from
equation~\eqref{eq:ch09:ncfsk}.
\item Derive the Rayleigh-fading BER of BPSK~\eqref{eq:ch09:rayleigh}. How much average SNR is needed
for $10^{-4}$? With two-branch maximal-ratio combining (use the formula in the figure script or
Chapter~\ref{ch:mimo}), how much is needed per branch?
\item A 4096-QAM Wi-Fi~7 transmitter must meet $-38$\,dB EVM with 3\,dB of margin. The PA contributes
$-45$\,dB, IQ imbalance $-50$\,dB and the DAC chain $-47$\,dB. What is the largest rms phase noise, in
degrees, that the synthesiser may have?
\item Derive the max-log LLRs~\eqref{eq:ch09:llr16} for Gray-labelled 16-QAM. Then compute the exact
LLRs~\eqref{eq:ch09:llr} for the sample of the worked example and find where the max-log values
are most in error.
\item Show that among all distributions on a fixed set of points with a given average energy, the
Maxwell--Boltzmann distribution~\eqref{eq:ch09:mb} maximises the entropy (use a Lagrange
multiplier). Why does a larger $\lambda$ lower the average energy?
\item \emph{(Simulation)} Using \texttt{gmsk\_baseband} from the chapter's figure script, generate MSK
and GMSK with $BT=0.5$, $0.3$ and $0.2$. Measure the 99\% power bandwidth of each, and the eye opening
of the in-phase component after a filter matched to the first Laurent pulse. Plot the trade-off.
\item \emph{(Simulation)} Implement the Maxwell--Boltzmann-shaped 64-QAM of
Figure~\ref{fig:ch09_shaping}. By Monte Carlo, estimate the mutual information at $E_s/N_0=15$\,dB
for $\lambda$ from 0 to 0.1 (on the $\pm1,\dots,\pm7$ grid), find the best $\lambda$, and compare the
shaping gain with the figure. Then add a MAP demapper with the correct priors and measure the gain
in pre-decoding bit error rate over an ML demapper that ignores them.
\end{enumerate}

\furtherreading
\begin{itemize}
\item J.~G.~Proakis and M.~Salehi, \emph{Digital Communications}, 5th ed., McGraw-Hill, 2008. The
standard graduate reference; Chapters 3 and 4 cover signal space, all the classical modulations and
their optimal detectors in depth.
\item J.~M.~Wozencraft and I.~M.~Jacobs, \emph{Principles of Communication Engineering}, Wiley, 1965.
The book that made signal space the language of the field; its geometric treatment of detection is
still a model of clarity.
\item M.~K.~Simon, S.~M.~Hinedi and W.~C.~Lindsey, \emph{Digital Communication Techniques: Signal
Design and Detection}, Prentice Hall, 1995. Exhaustive on coherent, differential and noncoherent
detection, including multiple-symbol differential detection.
\item M.~K.~Simon and M.-S.~Alouini, \emph{Digital Communication over Fading Channels}, 2nd ed.,
Wiley, 2005. The reference for error probabilities in fading, built on Craig's formula and its
generalisations.
\item J.~B.~Anderson, T.~Aulin and C.-E.~Sundberg, \emph{Digital Phase Modulation}, Plenum, 1986.
The definitive treatment of continuous-phase modulation.
\item S.~Pasupathy, ``Minimum shift keying: a spectrally efficient modulation,'' \emph{IEEE
Communications Magazine}, 17(4), 1979; and K.~Murota and K.~Hirade, ``GMSK modulation for digital
mobile radio telephony,'' \emph{IEEE Trans. Communications}, 29(7), 1981. The two papers behind MSK
and GMSK in practice.
\item G.~Ungerboeck, ``Channel coding with multilevel/phase signals,'' \emph{IEEE Trans. Information
Theory}, 28(1), 1982; and G.~D.~Forney Jr.\ and G.~Ungerboeck, ``Modulation and coding for linear
Gaussian channels,'' \emph{IEEE Trans. Information Theory}, 44(6), 1998, a masterly survey of
coding, shaping and the SNR gap.
\item G.~Caire, G.~Taricco and E.~Biglieri, ``Bit-interleaved coded modulation,'' \emph{IEEE Trans.
Information Theory}, 44(3), 1998. The analysis behind the architecture of every modern standard.
\item G.~B\"ocherer, F.~Steiner and P.~Schulte, ``Bandwidth efficient and rate-matched low-density
parity-check coded modulation,'' \emph{IEEE Trans. Communications}, 63(12), 2015. Probabilistic
amplitude shaping; see also F.~R.~Kschischang and S.~Pasupathy, ``Optimal nonuniform signaling for
Gaussian channels,'' \emph{IEEE Trans. Information Theory}, 39(3), 1993.
\item Standards: ETSI EN~302~307-1 and -2 (DVB-S2 and S2X constellations and ring ratios); IEEE
802.11ax-2021 and 802.11be (Wi-Fi MCS tables and EVM limits); 3GPP TS~38.211 (NR modulation mapping)
and TS~38.104 (base-station EVM).
\end{itemize}
